\loai{Loại 1: Chuyển động Brown}
\begin{ex}
	Thí nghiệm của Brown quan sát
	\choice
	{hạt phấn hoa bằng kính hiển vi}
	{\True chuyển động của hạt phấn hoa trong nước bằng kính hiển vi}
	{cánh hoa trong nước bằng kính hiển vi}
	{chuyển động của cánh hoa}
	\loigiai{
		\begin{itemize}
			\item Thí nghiệm của Brown quan sát chuyển động của các hạt phấn hoa trong nước, thấy các hạt phấn chuyển động hỗn độn, không ngừng
		\end{itemize}
	}
	\haidong{3}
\end{ex}

%\begin{ex}
%	Chuyển động Brown là chuyển động hỗn loạn không ngừng của
%	\choice
%	{các chất}
%	{các phân tử}
%	{các nguyên tử}
%	{\True các hạt rất nhỏ}
%	\loigiai{
%		\begin{itemize}
%			\item Thí nghiệm của Brown được coi là một trong những thí nghiệm chứng tỏ các phân tử chuyển động hỗn loạn, không ngừng.
%		\end{itemize}
%	}
%	\haidong{3}\end{ex}
\begin{ex}
	Trong thí nghiệm của Brown, các hạt phấn hoa chuyển động hỗn độn không ngừng vì
	
	\choice
	{giữa chúng có khoảng cách}
	{chúng là các phân tử}
	{\True các phân tử nước chuyển động không ngừng, va chạm vào chúng từ mọi phía}
	{chúng là các thực thể sống}
	\loigiai{
		\begin{itemize}
			\item Trong thí nghiệm của Brown, sở dĩ các hạt phấn hoa lại chuyển động hỗn độn không ngừng là vì các phân tử nước chuyển động không ngừng va chạm vào các hạt phấn hoa từ mọi phía.
		\end{itemize}
	}
	\haidong{3}\end{ex}


\begin{ex} \nguon{1}
	immini[thm]{Chuyển động hỗn loạn, không ngừng của các hạt rất nhỏ dưới tác dụng của các phân tử chất lỏng hoặc chất khí được gọi là chuyển động Brown}
	{\includegraphics[width=3cm]{images/2024_05_26_010fc85c05a422f9b1ffg-01(1)}}
\choiceTF[t]
{Chuyển động Brown do một nhà Vật Lí học phát hiện ra}
{Khi quan sát các hạt phấn hoa, Brown khẳng định cấu tạo chất bao gồm các phân tử}
{\True Nguyên nhân của chuyển động Brown là do các phân tử nước luôn luôn chuyển động hỗn loạn, va chạm vào các hạt phấn hoa từ nhiều phía}
{Hạt phấn hoa chuyển động có quy luật xác định}	
	\loigiai{
		\begin{itemchoice}
			\itemch 
			\itemch 
			\itemch 
			\itemch 
		\end{itemchoice}
%				\item SAI SAI ĐÚNG SAI
				Chuyển động Brown được phát hiện khi quan sát các hạt phấn hoa chuyển động hỗn loạn không có quy luật ở trong nước.
				Nguyên nhân là do các phân tử nước luôn luôn chuyển động hỗn loạn, va chạm vào các hạt phấn hoa từ nhiều phía.
	}
\end{ex}

%\begin{ex} \nguon{1}
%	Chuyển động Brown được phát hiện thế nào? Nguyên nhân là do đâu?
%	\loigiai{
%		immini[thm]{\begin{itemize}
%				\item Chuyển động Brown được phát hiện khi quan sát các hạt phấn hoa chuyển động hỗn loạn không có quy luật ở trong nước.
%				\item Nguyên nhân là do các phân tử nước luôn luôn chuyển động hỗn loạn, va chạm vào các hạt phấn hoa từ nhiều phía.
%		\end{itemize}}
%		{\includegraphics[width=3cm]{images/2024_05_26_010fc85c05a422f9b1ffg-01(1)}}
%	}
%\end{ex}


\loai{Loại 2: Cấu tạo chất}
\begin{ex}
	Phát biểu nào sau đây là \textbf{sai} khi nói về cấu tạo chất?	
	\choice
	{Các phân tử luôn luôn chuyển động không ngừng}
	{Các phân tử chuyển động càng nhanh thì nhiệt độ của vật càng cao}
	{\True Các phân tử luôn luôn đứng yên và chỉ chuyển động khi nhiệt độ của vật cao}
	{Các chất được cấu tạo từ các hạt riêng biệt gọi là phân tử}
	\loigiai{
		\begin{itemize}
			\item Các phân tử chuyển động không ngừng.
			Chuyển động của các phân tử được gọi là chuyển động nhiệt.
		\end{itemize}
	}
	\haidong{2}\end{ex}

\begin{ex}
	Nguyên tử, phân tử \textbf{không} có tính chất nào sau đây?
	\choice
	{Chuyển động không ngừng}
	{Giữa chúng có khoảng cách}
	{\True Nở ra khi nhiệt độ tăng, co lại khi nhiệt độ giảm}
	{Chuyển động càng nhanh khi nhiệt độ càng cao}
	\loigiai{
		\begin{itemize}
			\item Nguyên tử, phân tử không có tính chất nở ra khi nhiệt độ tăng, co lại khi nhiệt độ giảm, chỉ có chuyển động càng nhanh khi nhiệt độ càng cao
		\end{itemize}
	}
	\haidong{2}\end{ex}


\loai{Loại 3: Mô hình động học phân tử}
\begin{ex}
	Theo mô hình động học phân tử thì các chất
	\choice
	{\True được cấu tạo từ các hạt riêng biệt vô cùng nhỏ bé gọi là các phân tử}
	{liền một khối}
	{không được cấu tạo từ các hạt riêng biệt}
	{liền một khối, không được cấu tạo từ các hạt riêng biệt}
	\loigiai{
		\begin{itemize}
			\item Các chất được cấu tạo từ các hạt riêng biệt là các phân tử.
		\end{itemize}
		
	}
	\haidong{2}\end{ex}

\begin{ex}
	Phát biểu nào sau đây \textbf{không} đúng với mô hình động học phân tử?
	
	\choice
	{Các chất được cấu tạo từ các hạt riêng biệt là phân tử}
	{Các phân tử chuyển động không ngừng}
	{\True Tốc độ chuyển động của các phân tử cấu tạo nên vật càng lớn thì thể tích của vật càng lớn}
	{Giữa các phân tử có lực tương tác gọi là lực liên kết phân tử}
	\loigiai{
		\begin{itemize}
			\item Các phân tử chuyển động càng nhanh thì nhiệt độ của vật do chúng tạo nên càng cao
		\end{itemize}
		
	}
	\haidong{2}\end{ex}




\newpage 
\begin{ex} \nguon{4}
	Phát biểu nào sau đây là đúng khi nói về mô hình động học phân tử?
	\choice
	{Các phân tử cấu tạo nên các chất luôn đứng yên tại một vị trí cố định}
	{\True Các phân tử cấu tạo nên các chất luôn chuyển động không ngừng}
	{Các phân tử cấu tạo nên các chất có thể chuyển động hoặc đứng yên tùy vào đó là chất rắn, chất lỏng hay chất khí}
	{Các phân tử cấu tạo nên chất lỏng và chất khí thì chuyển động, các hạt cấu tạo nên chất rắn thì đứng yên}
	\loigiai{
		\begin{itemize}
			\item Theo mô hình động học phân tử, các phân tử cấu tạo nên các chất luôn chuyển động không ngừng.
			\item Chuyển động của các phân tử xảy ra ngay cả trong trạng thái rắn, lỏng và khí, mặc dù theo những cách khác nhau.
		\end{itemize}
	}
\end{ex}

\begin{ex} \nguon{4}
	Theo mô hình động học phân tử, nhiệt độ của vật cao hay thấp là do
	\choice
	{\True sự chuyển động nhanh hay chậm của các phân tử cấu tạo nên vật}
	{số phân tử cấu tạo nên vật là nhiều hay ít}
	{mật độ các phân tử trên một đơn vị thể tích là lớn hay nhỏ}
	{khối lượng của các phân tử là nặng hay nhẹ}
	\loigiai{
		\begin{itemize}
			\item Nhiệt độ của vật được xác định bởi sự chuyển động nhanh hay chậm của các phân tử cấu tạo nên vật.
			\item Chuyển động của các phân tử nhanh hơn tương ứng với nhiệt độ cao hơn, trong khi chuyển động chậm hơn tương ứng với nhiệt độ thấp hơn.
		\end{itemize}
	}
\end{ex}

\begin{ex} \nguon{4}
	Mô hình động học phân tử được xây dựng dựa trên quan điểm nào về cấu tạo chất?
	\choice
	{Dựa trên quan điểm chất có cấu tạo liên tục}
	{\True Dựa trên quan điểm chất có cấu tạo gián đoạn}
	{Dựa trên cả hai quan điểm là chất có cấu tạo liên tục và gián đoạn}
	{Dựa trên quan điểm về sự bền vững của các phân tử}
	\loigiai{
		\begin{itemize}
			\item Mô hình động học phân tử được xây dựng dựa trên quan điểm chất có cấu tạo gián đoạn.
			\item Theo quan điểm này, chất được cấu tạo từ các phân tử rời rạc, không liên tục và luôn chuyển động.
		\end{itemize}
	}
\end{ex}


\begin{ex}
	Tính chất nào sau đây \textbf{không} phải là của phân tử?
	\choice
	{Chuyển động không ngừng}
	{Giữa các phân tử có khoảng cách}
	{\True Có lúc đứng yên, có lúc chuyển động}
	{Chuyển động càng nhanh thì nhiệt độ của vật càng cao}
	\loigiai{
		\begin{itemize}
			\item Các nguyên tử, phân tử chuyển động không ngừng
		\end{itemize}
	}
	\haidong{2}\end{ex}
\begin{ex} \nguon{1}
	Phát biểu nào sau đây \textbf{không} đúng?
	\choice
	{Các chất được cấu tạo từ các hạt gọi là phân tử}
	{Lực tương tác giữa các phân tử ở thể rắn lớn hơn lực tương tác giữa các phân tử ở thể lỏng và thể khí}
	{Các phân tử chất lỏng dao động xung quanh các vị trí cân bằng không cố định}
	{\True Các phân tử đứng sát nhau và giữa chúng không có khoảng cách}
	\loigiai{
		\begin{itemize}
			\item Phát biểu "Các phân tử đứng sát nhau và giữa chúng không có khoảng cách" là không đúng. Trong thực tế, ngay cả khi các phân tử ở gần nhau nhất, giữa chúng vẫn tồn tại khoảng cách nhất định do sự tồn tại của các lực tương tác và các quy luật cơ học lượng tử. 
			\item Những phát biểu còn lại đều đúng:
			\begin{itemize}
				\item "Các chất được cấu tạo từ các hạt gọi là phân tử" là đúng vì theo mô hình cấu tạo của vật chất, tất cả các chất đều được cấu thành từ các phân tử.
				\item "Lực tương tác giữa các phân tử ở thể rắn lớn hơn lực tương tác giữa các phân tử ở thể lỏng và thể khí" là đúng vì ở thể rắn, các phân tử sắp xếp gần nhau và có lực hút lớn, tạo thành cấu trúc cố định, trong khi ở thể lỏng và thể khí, các lực tương tác yếu hơn.
				\item "Các phân tử chất lỏng dao động xung quanh các vị trí cân bằng không cố định" cũng là đúng vì trong chất lỏng, các phân tử có vị trí cân bằng thay đổi liên tục do chúng có thể trượt qua nhau nhưng không tách rời hoàn toàn.
			\end{itemize}
		\end{itemize}
	}
\end{ex}









\loai{Loại 4: Lực tương tác phân tử}
\begin{ex}
	Phát biểu nào sau đây là đúng khi nói về lực tương tác giữa các phân tử?
	\choice
	{\True Giữa các phân tử đồng thời có cả lực hút và lực đẩy}
	{Giữa các phân tử chỉ có lực hút hoặc lực đẩy}
	{Giữa các phân tử chỉ có lực đẩy}
	{Giữa các phân tử chỉ có lực hút}
	\loigiai{
		\begin{itemize}
			\item Giữa các phân tử có lực hút và lực đẩy, gọi chung là lực liên kết phân tử.
		\end{itemize}
	}
	\haidong{2}\end{ex}

\begin{ex}
	Lực liên kết giữa các phân tử	
	\choice
	{là lực hút}
	{là lực đẩy}
	{tuỳ thuộc vào thể của nó, ở thể rắn là lực hút còn ở thể khí là lực đẩy}
	{\True gồm cả lực hút và lực đẩy}
	\loigiai{
		\begin{itemize}
			\item Giữa các phân tử có lực hút và lực đẩy, gọi chung là lực liên kết phân tử.
		\end{itemize}
	}
	\haidong{2}\end{ex}



\begin{ex}%[1P5B1-1][Loai1]
	Phát biểu nào sau đây nói về lực tương tác phân tử là \textbf{không} đúng?
	\choice
	{Lực tương tác phân tử đáng kể khi các phân tử ở rất gần nhau}
	{Lực hút phân tử có thể lớn hơn lực đẩy phân tử}
	{\True Lực hút phân tử không thể lớn hơn lực đẩy phân tử}
	{Lực hút phân tử có thể bằng lực đẩy phân tử}
	\loigiai{
		\begin{itemize}
			\item Giữa các phân tử luôn có cả lực hút và lực đẩy. Lực hút và sẽ lớn hơn hay nhỏ hơn lực đẩy tùy vào khoảng cách giữa các phân tử.
	\end{itemize}} 
	%<MyLT3>
	\haidong{2}\end{ex}
\newpage 
\begin{ex} %[1P5Y1-1][Loai1]
	Khi khoảng cách giữa các phân tử rất nhỏ so với kích thước phân tử, thì giữa các phân tử
	\choice
	{có cả lực hút và lực đẩy, nhưng lực đẩy nhỏ hơn lực hút}
	{\True có cả lực hút và lực đẩy, nhưng lực đẩy lớn hơn lực hút}
	{chỉ có lực đẩy}
	{chỉ có lực hút}	
	\loigiai{
		\begin{itemize}
			\item Khi khoảng cách giữa các phân tử rất nhỏ, thì giữa các phân tử có cả lực hút và lực đẩy, nhưng lực đẩy lớn hơn lực hút.
		\end{itemize}
	}%<MyLT1>
	\haidong{2}\end{ex}



%\begin{ex} \nguon{4}
%	Phát biểu nào sau đây về chuyển động của phân tử là \textbf{không} đúng?
%	\choice
%	{\True Chuyển động của phân tử là do lực tương tác phân tử gây ra}
%	{Các phân tử chuyển động không ngừng}
%	{Các phân tử chuyển động chậm thì nhiệt độ của vật thấp}
%	{Các phân tử khí chuyển động theo đường thẳng giữa hai lần va chạm}
%	\loigiai{
%		\begin{itemize}
%			\item Chuyển động của phân tử không phải là do lực tương tác phân tử gây ra. Thực chất, các phân tử chuyển động do năng lượng nhiệt mà chúng có.
%			\item Các phân tử chuyển động không ngừng và nhiệt độ của vật liên quan đến tốc độ chuyển động của các phân tử.
%			\item Các phân tử khí chuyển động theo đường thẳng giữa hai lần va chạm.
%		\end{itemize}
%	}
%\end{ex}
\loai{Loại 5. Chuyển động nhiệt của phân tử}

\begin{ex}
	Phát biểu nào sau đây nói về chuyển động của phân tử là \textbf{không} đúng?
	\choice
	{\True Chuyển động của phân tử hoàn toàn là do lực tương tác phân tử gây ra}
	{Các phân tử chuyển động không ngừng}
	{Các phân tử chuyển động càng nhanh thì nhiệt độ của vật càng cao}
	{Các phân tử khí chuyển hỗn loạn không ngừng về mọi hướng}
	\loigiai{
		\begin{itemize}
			\item Giữa các phân tử có lực hút và lực đẩy, gọi chung là lực liên kết phân tử.
			\item Lực tương tác phân tử có thể gây thay đổi chuyển động chứ không gây ra chuyển động.
		\end{itemize}
	}
	\haidong{2}
\end{ex}





\begin{ex} \nguon{1}
	Vận tốc chuyển động nhiệt của các phân tử có ảnh hưởng trực tiếp đến đại lượng nào sau đây?
	\choice
	{\True Nhiệt độ của vật}
	{Trọng lượng riêng của vật}
	{Khối lượng của vật}
	{Thể tích của vật}
	\loigiai{
		\begin{itemize}
			\item Vận tốc chuyển động của các phân tử trong một chất có liên quan mật thiết đến nhiệt độ của chất đó. Khi nhiệt độ tăng, năng lượng nhiệt của các phân tử tăng lên, khiến chúng chuyển động nhanh hơn. Ngược lại, khi nhiệt độ giảm, vận tốc chuyển động của các phân tử cũng giảm theo.		
			\item Các đại lượng khác như trọng lượng riêng, khối lượng, và thể tích của vật không ảnh hưởng trực tiếp đến vận tốc chuyển động của các phân tử. Do đó, đáp án đúng là nhiệt độ của vật.
		\end{itemize}
	}
\end{ex}

\begin{ex} \nguon{1}
	Phát biểu nào sau đây là đúng?
	\choice
	{Ở cùng nhiệt độ, các phân tử của các chất khác nhau sẽ có vận tốc chuyển động giống nhau}
	{Các nguyên tử, phân tử cấu tạo nên vật không chuyển động}
	{\True Nhiệt độ càng cao thì các nguyên tử, phân tử chuyển động càng nhanh}
	{Các vật được cấu tạo liền một khối}
	\loigiai{
		\begin{itemize}
			\item Phát biểu đúng là "Nhiệt độ càng cao thì các nguyên tử, phân tử chuyển động càng nhanh". Điều này dựa trên nguyên lý cơ bản của thuyết động học phân tử, theo đó khi nhiệt độ tăng, các phân tử sẽ chuyển động nhanh hơn.		
			\item Phát biểu “Ở cùng nhiệt độ, các phân tử của các chất khác nhau sẽ có vận tốc chuyển động giống nhau” không hoàn toàn chính xác, vì vận tốc trung bình của phân tử phụ thuộc chủ yếu vào khối lượng phân tử và nhiệt độ, không chỉ riêng nhiệt độ.		
			\item Phát biểu “Các nguyên tử, phân tử cấu tạo nên vật không chuyển động” là sai, vì trong thực tế, các nguyên tử và phân tử luôn có chuyển động nhiệt.		
			\item Phát biểu “Các vật được cấu tạo liền một khối” cũng là không đúng, vì các vật thực tế được tạo thành từ các nguyên tử và phân tử riêng biệt.
			\item Do đó, đáp án chính xác là “Nhiệt độ càng cao thì các nguyên tử, phân tử chuyển động càng nhanh”.
		\end{itemize}
	}
\end{ex}



%\begin{ex} \nguon{4}
%	Khi các phân tử nhận được thêm năng lượng thì các phân tử chuyển động hỗn loạn
%	\choice
%	{càng nhanh, khoảng cách trung bình giữa các phân tử càng xa}
%	{càng chậm, khoảng cách trung bình giữa các phân tử càng gần}
%	{\True càng nhanh, khoảng cách trung bình giữa các phân tử càng xa}
%	{càng nhanh, khoảng cách trung bình giữa các phân tử càng gần}
%	\loigiai{
%		Khi phân tử nhận thêm năng lượng, chúng chuyển động nhanh hơn và khoảng cách trung bình giữa chúng tăng lên.
%	}
%\end{ex}




\loai{Loại 6: Xác định các thể}


\begin{ex}
	Phát biểu nào sau đây là \textbf{sai}? Cùng một khối lượng của một chất nhưng khi ở các thể khác nhau thì sẽ khác nhau về
	\choice
	{thể tích}
	{khối lượng riêng}
	{\True kích thước của các nguyên tử}
	{trật tự của các nguyên tử}
	
	\loigiai{
		\begin{itemize}
			\item Sử dụng lý thuyết các thể của chất: Kích thước của các nguyên tử ở các thể khác nhau là như nhau
		\end{itemize}
	}
	\haidong{1}\end{ex}

\begin{ex}
	Các chất có thể tồn tại ở những thể cơ bản nào?
	\choice
	{Thể rắn}
	{Thể lỏng}
	{Thể khí}
	{\True Thể rắn, thể lỏng, thể khí}
	\loigiai{
		\begin{itemize}
			\item Ba thể cơ bản của các chất: Thể rắn, thể lỏng, thể khí
		\end{itemize}
	}
	\haidong{3}\end{ex}

\begin{ex}
	Chất nào sau đây có thể tồn tại ở cả ba thể rắn, lỏng, khí ở điều kiện tự nhiên trên Trái Đất?
	
	\choice
	{\True Nước}
	{Oxygen}
	{Nitrogen}
	{Carbon dioxide}
	\loigiai{
		
		\begin{itemize}
			\item Nước có thể tồn tại ở thể rắn (nước đá, băng, tuyết), thể lỏng (nước lỏng), thể khí (hơi nước).
		\end{itemize}
		
	}
	\haidong{2}\end{ex}

\begin{ex}
	Cho các chất sau: đồng (copper); nitrogen; vàng (gold); nước; oxygen; dầu hỏa; sắt (iron). Số chất ở thể rắn trong điều kiện thường là
	
	\choice
	{1}
	{2}
	{\True 3}
	{4}
	\loigiai{
		
		\begin{itemize}
			\item Các chất ở thể rắn trong điều kiện thường là: đồng (copper); vàng (gold); sắt (iron).
		\end{itemize}
		
	}
	\haidong{2}\end{ex}

\begin{ex}
	Chất nào sau đây ở thể khí khi ở điều kiện thường?
	
	\choice
	{Muối ăn}
	{Đường}
	{Giấm}
	{\True Carbon dioxide}
	\loigiai{
		
		\begin{itemize}
			\item 	Carbon dioxide là chất khí ở điều kiện thường.
		\end{itemize}
		
	}
	\haidong{2}\end{ex}








\loai{Loại 7: Cấu trúc chất rắn}
\begin{ex}
	Hình dưới đây biểu diễn mô hình cấu tạo phân tử của ba chất $A,\ B$ và $C$.
	\begin{center}
		\begin{tikzpicture}[scale=1.1,thick,>=stealth']
			\tkzDefPoints{0.2/2/a, .8/2/b, 3.2/0/c, 3.5/.45/d}
			\draw (-0.5,-0.5) -- (4,-0.5) -- (4,2.5) -- (-0.5,2.5) -- (-0.5,-0.5);
			\node [draw, thick, shape=circle, minimum width=.65cm, minimum height=.5cm, anchor=center] at (b) {};
			\node [draw, thick,fill=white, shape=circle, minimum width=.65cm, minimum height=.5cm, anchor=center] at (a) {};
			\node [draw, thick,shape=circle, minimum width=.65cm, minimum height=.5cm, anchor=center] at (d) {};
			\node [draw, thick,fill=white, shape=circle, minimum width=.65cm, minimum height=.5cm, anchor=center] at (c) {};
			\node at (1.75,-1) {$A$};
			
			\begin{scope}[shift={(5,0)},rotate=0]
				\tkzDefPoints{0.5/2/a, 1.5/2/b, 2.5/2/c, 3.5/2/d}
				\tkzDefPoints{0.5/1/x, 1.5/1/y, 2.5/1/z, 3.5/1/o}
				\tkzDefPoints{0.5/0/k, 1.5/0/m, 2.5/0/n, 3.5/0/p}
				
				\draw (0,-0.5) -- (4,-0.5) -- (4,2.5) -- (0,2.5) -- (0,-0.5);
				\node [draw, thick, shape=circle, minimum width=.95cm, minimum height=.95cm, anchor=center] at (a) {};
				\node [draw, thick, shape=circle, minimum width=.95cm, minimum height=.95cm, anchor=center] at (b) {};
				\node [draw, thick, shape=circle, minimum width=.95cm, minimum height=.95cm, anchor=center] at (c) {};
				\node [draw, thick, shape=circle, minimum width=.95cm, minimum height=.95cm, anchor=center] at (d) {};
				
				\node [draw, thick, shape=circle, minimum width=.95cm, minimum height=.95cm, anchor=center] at (x) {};
				\node [draw, thick, shape=circle, minimum width=.95cm, minimum height=.95cm, anchor=center] at (y) {};
				\node [draw, thick, shape=circle, minimum width=.95cm, minimum height=.95cm, anchor=center] at (z) {};
				\node [draw, thick, shape=circle, minimum width=.95cm, minimum height=.95cm, anchor=center] at (o) {};
				
				\node [draw, thick, shape=circle, minimum width=.95cm, minimum height=.95cm, anchor=center] at (k) {};
				\node [draw, thick, shape=circle, minimum width=.95cm, minimum height=.95cm, anchor=center] at (m) {};
				\node [draw, thick, shape=circle, minimum width=.95cm, minimum height=.95cm, anchor=center] at (n) {};
				\node [draw, thick, shape=circle, minimum width=.95cm, minimum height=.95cm, anchor=center] at (p) {};
				
				\node at (2,-1) {$B$};
			\end{scope}
			
			\begin{scope}[shift={(10,0)},rotate=0]
				\tkzDefPoints{0.5/2/a, 1.25/1.65/b, 2.5/1.3/c, 2.3/1.75/d}
				\tkzDefPoints{0.75/1.5/x, 3.5/2/y, 4.15/1.5/z, 3.65/1.5/o}
				\tkzDefPoints{1.5/.5/k, 1/0.25/m, 2.75/0/n, 3.5/0.5/p, 3/0.45/q}
				
				\draw (0,-0.5) -- (4.5,-0.5) -- (4.5,2.5) -- (0,2.5) -- (0,-0.5);
				\node [draw, thick, shape=circle, minimum width=.6cm, minimum height=.5cm, anchor=center] at (b) {};
				\node [draw, thick,fill=white, shape=circle, minimum width=.6cm, minimum height=.5cm, anchor=center] at (a) {};
				\node [draw, thick,fill=white, shape=circle, minimum width=.8cm, minimum height=.5cm, anchor=center] at (x) {};
				
				\node [draw, thick, shape=circle, minimum width=.6cm, minimum height=.5cm, anchor=center] at (c) {};
				\node [draw, thick,fill=blue,shape=circle, minimum width=.9cm, minimum height=.5cm, anchor=center] at (d) {};
				
				\node [draw, thick, shape=circle, minimum width=.6cm, minimum height=.5cm, anchor=center] at (y) {};
				\node [draw, thick,fill=white, shape=circle, minimum width=.6cm, minimum height=.5cm, anchor=center] at (z) {};
				\node [draw, thick,fill=white, shape=circle, minimum width=.8cm, minimum height=.5cm, anchor=center] at (o) {};
				
				\node [draw, thick, shape=circle, minimum width=.6cm, minimum height=.5cm, anchor=center] at (k) {};
				\node [draw, thick,fill=blue,shape=circle, minimum width=.9cm, minimum height=.5cm, anchor=center] at (m) {};
				
				\node [draw, thick, shape=circle, minimum width=.6cm, minimum height=.5cm, anchor=center] at (n) {};
				\node [draw, thick,fill=white, shape=circle, minimum width=.6cm, minimum height=.5cm, anchor=center] at (p) {};
				\node [draw, thick,fill=white, shape=circle, minimum width=.8cm, minimum height=.5cm, anchor=center] at (q) {};
				\node at (2.25,-1) {$C$};
			\end{scope}
			
			\node at (7,-2) {};
		\end{tikzpicture}
	\end{center}
	\vspace*{-1.5 cm}
	Mô hình chất rắn tinh thể giống với mô hình nào nhất sau đây?
	\choice
	{A}
	{\True B}
	{C}
	{A và B}
	\loigiai{
		\begin{itemize}
			\item Từ ba mô hình đã cho, chất B là chất rắn vì các phân tử ở gần nhau nhất và được sắp xếp theo trật tự xác định.
		\end{itemize}
		
	}

\end{ex}
%
\begin{ex}
	Đặc điểm nào sau đây là của chất rắn?
	\choice
	{\True Lực tương tác giữa các phân tử rất mạnh}
	{Không có hình dạng xác định}
	{Lực tương tác giữa các phân tử rất yếu}
	{Không có thể tích riêng xác định}
	\loigiai{
		\begin{itemize}
			\item Trong chất rắn, các phân tử ở rất gần nhau nên lực tương tác giữa các phân tử chất rắn rất mạnh.
		\end{itemize}
	}
	\haidong{2}\end{ex}

\begin{ex}
	Vật rắn có hình dạng xác định vì phân tử cấu tạo nên vật rắn
	\choice
	{không chuyển động}
	{đứng sát nhau}
	{chuyển động với vận tốc nhỏ không đáng kể}
	{\True chuyển động quanh một vị trí cân bằng xác định}
	\loigiai{
		\begin{itemize}
			\item Các phân tử chất rắn được giữ ở các vị trí cân bằng và mỗi phân tử chỉ có thể dao động xung quanh vị trí cân bằng xác định này nên các chất ở thể rắn có thể tích và hình dạng xác định.
		\end{itemize}
	}
	\haidong{2}\end{ex}

\begin{ex}
	Tại sao người ta không nén chất rắn lại để tiết kiệm diện tích?
	\choice
	{\True Do chất rắn liên kết chặt chẽ với nhau nên rất khó nén}
	{Do chất rắn liên kết không chặt chẽ với nhau nên khó nén}
	{Do các hạt của chất rắn chuyển động tự do}
	{Do người ta không muốn nén chất rắn}
	\loigiai{
		\begin{itemize}
			\item Trong chất rắn, các phân tử ở rất gần nhau nên lực tương tác giữa các phân tử chất rắn chặt chẽ, do đó nó có thể tích xác định và không thể nén được.
		\end{itemize}
	}
	\haidong{2}\end{ex}

\begin{ex}
	Những đặc điểm nào sau đây là của chất rắn?
	\choice
	{Có khối lượng và hình dạng xác định, không có thể tích xác định}
	{Có khối lượng xác định, hình dạng và thể tích không xác định}
	{\True Có khối lượng, hình dạng và thể tích xác định}
	{Có khối lượng và thể tích xác định, hình dạng không xác định}
	\loigiai{
		\begin{itemize}
			\item Các phân tử chất rắn được giữ ở các vị trí cân bằng và mỗi phân tử chỉ có thể dao động xung quanh vị trí cân bằng xác định này nên các chất ở thể rắn có thể tích và hình dạng xác định.
		\end{itemize}
	}
	\haidong{2}\end{ex}

\begin{ex}
	Phát biểu nào dưới đây là \textbf{không} đúng với thể rắn?
	\choice
	{Các phân tử rất gần nhau (cỡ kích thước phân tử)}
	{Các phân tử sắp xếp có trật tự}
	{Các phân tử dao động quanh vị trí cân bằng cố định}
	{\True Các phân tử dao động quanh vị trí cân bằng không xác định}
	\loigiai{
		\begin{itemize}
			\item Các phân tử chất rắn được giữ ở các vị trí cân bằng và mỗi phân tử
			chỉ có thể dao động xung quanh vị trí cân bằng xác định này.
		\end{itemize}
	}
	\haidong{2}\end{ex}

%\begin{ex}%[1P5Y1-1][Loai2]
%	Vị trí của các nguyên tử, phân tử trong chất rắn có đặc điểm nào sau đây?
%	\choice
%	{\True Các nguyên tử, phân tử nằm ở những vị trí xác định và chỉ có thể dao động xung quanh các vị trí xác định này}
%	{Các nguyên tử, phân tử nằm ở những vị trí xác định và không dao động}
%	{Các nguyên tử, phân tử không có vị trí cố định mà luôn thay đổi}
%	{Các nguyên tử, phân tử nằm ở những vị trí cố định, sau một thời gian nào đó chúng lại chuyển sang một vị trí cố định khác}
%	\loigiai{
%		\begin{itemize}
%			\item Các nguyên tử, phân tử nằm ở những vị trí xác định và chỉ có thể dao động xung quanh các vị trí cân bằng xác định này.
%	\end{itemize}}%<MyLT2>
%	\haidong{2}\end{ex}

\begin{ex}
	Các nguyên tử trong một miếng sắt có tính chất nào sau đây?
	\choice
	{Khi nhiệt độ tăng thì nở ra}
	{Khi nhiệt độ giảm thì co lại}
	{\True Khoảng cách rất gần nhau}
	{Khoảng cách rất xa nhau}
	\loigiai{
		\begin{itemize}
			\item Các nguyên tử trong kim loại nói riêng và chất rắn nói chung thường đứng rất gần nhau.
		\end{itemize}		
	}
	\haidong{2}\end{ex}
\begin{ex} \nguon{1}
			Tại sao cùng là carbon nhưng than chì rất mềm còn kim cương lại vô cùng cứng?
			\begin{center}
					 includegraphics[width=8cm]{images/80316}				
			\end{center}
\choiceTF[t]
{Do than chì và kim cương được cấu tạo từ hai chất khác nhau}
{Do than chì là chất rắn vô định hình}
{\True Do than chì có cấu trúc tinh thể khác kim cương}
{\True Do cấu trúc tinh thể của than chì xếp theo lớp nên các nguyên tử carbon liên kết với nhau yếu}	
	\loigiai{
		\begin{itemchoice}
		\itemch Than chì và kim cương đều được cấu tạo từ các nguyên tử carbon nhưng do cấu trúc tinh thể khác nhau nên tính chất vật lí của chúng khác nhau.
		\itemch 
		\itemch Cấu trúc tinh thể kim cương là lập phương vô cùng bền vững. Còn cấu trúc than chì là cấu trúc xếp theo lớp 2D có thể trượt lên nhau, chúng liên kết với nhau rất yếu và tách khỏi nhau dễ dàng.
		\itemch 
	\end{itemchoice}
} 
\end{ex}


\begin{ex}
	Hình 1 mô tả cấu trúc của một chất rắn. Trong hình 2, hình nào thể hiện đúng nhất cấu trúc của chất rắn khi bị nung nóng?
	\begin{center}
		\begin{tikzpicture}[scale=1,thick,>=stealth']
		\tkzDefPoints{0.5/2/a, 1.5/2/b, 2.5/2/c, 3.5/2/d}
		\tkzDefPoints{0.5/1/x, 1.5/1/y, 2.5/1/z, 3.5/1/o}
		\tkzDefPoints{0.5/0/k, 1.5/0/m, 2.5/0/n, 3.5/0/p}	
		\draw (0,-0.5) -- (2,-0.5) -- (2,2.5) -- (0,2.5) -- (0,-0.5);
		\shade[ball color=blue] (a) circle (2ex);
		\shade[ball color=blue] (b) circle (2ex);
		\shade[ball color=blue] (x) circle (2ex);
		\shade[ball color=blue] (y) circle (2ex);
		\shade[ball color=blue] (k) circle (2ex);
		\shade[ball color=blue] (m) circle (2ex);
		\node at (1,-1) {Hình 1};
		
		\begin{scope}[shift={(4,0)},rotate=0]
			\tkzDefPoints{0.5/3/a, 1.5/3/b, 3.5/3/c, 4.5/3/d}
			\tkzDefPoints{0.5/1/x, 1.5/1/y, 3.5/1/z, 4.5/1/o}
			\tkzDefPoints{0.5/2/k, 1.5/2/m, 3.5/2/n, 4.5/2/p}
			\tkzDefPoints{6.5/3/a', 8/3/b', 6.5/1/c', 8/1/d', 6.5/2/k', 8/2/m'}		
			\draw (0,-0.5) -- (9,-0.5) -- (9,4.5) -- (0,4.5) -- (0,-0.5);
			\shade[ball color=blue] (a) circle (1.75ex);
			\shade[ball color=blue] (b) circle (1.75ex);
			\shade[ball color=blue] (x) circle (1.75ex);
			\shade[ball color=blue] (y) circle (1.75ex);
			\shade[ball color=blue] (k) circle (1.75ex);
			\shade[ball color=blue] (m) circle (1.75ex);
			
			\shade[ball color=blue] (c) circle (2.5ex);
			\shade[ball color=blue] (d) circle (2.5ex);
			\shade[ball color=blue] (z) circle (2.5ex);
			\shade[ball color=blue] (o) circle (2.5ex);
			\shade[ball color=blue] (n) circle (2.5ex);
			\shade[ball color=blue] (p) circle (2.5ex);
			
			\shade[ball color=blue] (a') circle (2ex);
			\shade[ball color=blue] (b') circle (2ex);
			\shade[ball color=blue] (c') circle (2ex);
			\shade[ball color=blue] (d') circle (2ex);
			\shade[ball color=blue] (k') circle (2ex);
			\shade[ball color=blue] (m') circle (2ex);
			\node at (4,-1) {Hình 2};
			\node at (1,0) {A};
			\node at (4,0) {B};
			\node at (7,0) {C};
		\end{scope}
		
		
	\end{tikzpicture}
	\end{center}
	\choice
	{Hình A}
	{Hình B}
	{\True Hình C}
	{Hình A và C}
	\loigiai{
		\begin{itemize}
			\item Hình thể hiện đúng nhất cấu trúc của chất rắn khi bị nung nóng là: \textbf{C}.
			\item 		\textbf{Lý do:} Khi chất rắn bị nung nóng, các phân tử bắt đầu dao động mạnh hơn và khoảng cách giữa chúng tăng lên, làm cho mạng lưới trở nên kém trật tự hơn so với trạng thái nguyên thủy.
		\end{itemize}
	}
\end{ex}





\loai{Loại 8: Chất rắn kết tinh}


\begin{ex}
	Các nguyên tử sắp xếp theo một trật tự nhất định trong không gian được gọi là
	\choice
	{mạng nguyên tử}
	{mạng internet}
	{\True mạng tinh thể}
	{mạng nhện}
	\loigiai{
		\begin{itemize}
			\item Các phân tử liên kết chặt với nhau và sắp xếp theo một trật tự hình học xác
			định, tuần hoàn trong không gian, gọi là mạng tinh thể.
		\end{itemize}
	}
	\haidong{2}\end{ex}

\begin{ex}
	Trong các chất sau, chất nào \textbf{không} phải là chất rắn kết tinh?
	\choice
	{Muối ăn}
	{\True Thuỷ tinh}
	{Kim cương}
	{Thạch anh}
	\loigiai{
		\begin{itemize}
			\item Thủy tinh không có cấu trúc tinh thể, là chất rắn vô định hình
		\end{itemize}
	}
	\haidong{2}\end{ex}




\begin{ex}%[1P7B1-1][Loai1]
	Chất nào sau đây \textbf{không} có cấu trúc tinh thể?
	\choice
	{\True Miếng nhựa thông}
	{Hạt đường}
	{Viên kim cương}
	{Khối thạch anh}
	\loigiai{
		\begin{itemize}
			\item Miếng nhựa thông là chất rắn vô định hình, không có cấu trúc tinh thể
		\end{itemize}
	}	
	%<MyLT3>
	\haidong{2}\end{ex}




\begin{ex} \nguon{1}
	Trong tinh thể, các hạt (nguyên tử, phân tử, ion)
	\choice
	{\True dao động nhiệt xung quanh vị trí cân bằng}
	{đứng yên tại những vị trí xác định}
	{chuyển động hỗn độn không ngừng}
	{chuyển động trên quỹ đạo tròn xung quanh một vị trí xác định}
	\loigiai{
		\begin{itemize}
			\item Trong tinh thể, các nguyên tử, phân tử, hoặc ion không đứng yên mà dao động xung quanh các vị trí cân bằng của chúng. 
			\item Vị trí cân bằng này là các vị trí cấu trúc xác định mà các hạt chiếm giữ. Khi nhiệt độ tăng, dao động nhiệt của các hạt tăng lên, nhưng chúng vẫn dao động xung quanh vị trí cân bằng cố định. 
			\item Đây là lý do vì sao tinh thể có hình dạng và cấu trúc ổn định ở điều kiện bình thường.
		\end{itemize}
	}
\end{ex}

\begin{ex} \nguon{1}
	immini[thm]{Hình bên là cấu trúc tinh thể của
		\motcot
		{\True muối ăn}
		{kim cương}
		{thủy tinh}
		{đồng}}
	{\includegraphics[width=2cm]{images/2024_05_26_010fc85c05a422f9b1ffg-05}} 
	\loigiai{
		Muối ăn có cấu trúc lập phương.
	}
\end{ex}




\begin{ex} \nguon{4}
	Khi chất rắn kết tinh được nung nóng. Phát biểu nào sau đây là đúng?
	\choice
	{các phân tử vẫn dao động với biên độ không đổi, khoảng cách trung bình giữa các phân tử không đổi}
	{\True các phân tử dao động với biên độ tăng lên, khoảng cách trung bình giữa các phân tử tăng lên}
	{các phân tử dao động với biên độ không đổi, khoảng cách trung bình giữa các phân tử tăng lên}
	{các phân tử dao động với biên độ tăng lên, khoảng cách trung bình giữa các phân tử không đổi}
	\loigiai{
		Khi nung nóng, các phân tử trong chất rắn dao động mạnh hơn (biên độ tăng), làm cho khoảng cách giữa các phân tử cũng tăng lên.
	}
\end{ex}


\loai{Loại 9: Chất rắn vô định hình}
\begin{ex}
	Chất rắn vô định hình là chất rắn
	\choice
	{có cấu trúc tinh thể}
	{\True không có cấu trúc tinh thể}
	{có cấu trúc mạng lập phương}
	{có cấu trúc mạng lục giác}
	\loigiai{
		\begin{itemize}
			\item Chất rắn vô định hình: Không có cấu trúc tinh thể.
			\item Thuỷ tinh, nhựa đường, cao su,... là những chất rắn vô định hình.
		\end{itemize}
	}
	\haidong{2}\end{ex}

\begin{ex}
	Chất rắn nào sau đây thuộc dạng chất rắn vô định hình?
	\choice
	{Muối ăn}
	{Kim loại}
	{Thạch anh}
	{\True Nhựa đường}
	\loigiai{
		\begin{itemize}
			\item Chất rắn vô định hình: Không có cấu trúc tinh thể.
			\item Thuỷ tinh, nhựa đường, cao su,... là những chất rắn vô định hình.
		\end{itemize}
	}
	\haidong{3}\end{ex}

\begin{ex}
	Chất rắn nào dưới đây thuộc loại chất rắn vô định hình?
	\choice
	{Ống nano carbon}
	{\True Socola}
	{Nhôm}
	{Kim cương}
	\loigiai{
		\begin{itemize}
			\item Chất rắn vô định hình: Không có cấu trúc tinh thể.
			Socola là chất rắn vô định hình
		\end{itemize}	
	}
	\haidong{3}\end{ex}

\begin{ex}%[1P7B1-1][Loai2]
	Chất rắn nào dưới đây là chất rắn vô định hình?
	\choice
	{\True Thủy tinh}
	{Băng phiến}
	{Hợp kim}
	{Kim loại}
	\loigiai{
		\begin{itemize}
			\item Các chất rắn vô định hình như thủy tinh, các loại nhựa, cao su, ...
		\end{itemize}
	}%<MyLT1>
	\haidong{3}\end{ex}
\loai{Loại 10: Chất rắn kết tinh – Chất rắn vô định hình}


\begin{ex}
	Chất rắn nào dưới đây \textbf{không} phải là chất rắn vô định hình?
	\choice
	{\True Thạch anh}
	{Thủy tinh}
	{Sáp}
	{Cao su}
	\loigiai{
		\begin{itemize}
			\item Thạch anh có cấu trúc tinh thể, là chất rắn kết tinh
		\end{itemize}
	}
	\haidong{3}\end{ex}






\begin{ex} \nguon{1}
	Vật nào sau đây \textbf{không} có cấu trúc tinh thể?
	\choice
	{Hạt muối}
	{Viên kim cương}
	{Miếng thạch anh}
	{\True Chiếc cốc làm bằng thủy tinh}
	\loigiai{
		\begin{itemize}
			\item Trong số các vật được liệt kê, chiếc cốc làm bằng thủy tinh không có cấu trúc tinh thể. 
			\item Thủy tinh là một chất rắn vô định hình, nghĩa là các nguyên tử trong thủy tinh không sắp xếp theo một trật tự tuần hoàn và đều đặn như trong cấu trúc tinh thể. 
			\item Ngược lại, hạt muối, viên kim cương và miếng thạch anh đều có cấu trúc tinh thể, trong đó các nguyên tử hoặc ion được sắp xếp theo một mạng lưới đều đặn và lặp đi lặp lại.\end{itemize}
	}
\end{ex}




\loai{Loại 11: Cấu trúc chất lỏng}
\begin{ex}
	Đặc điểm nào sau đây là đặc điểm của thể lỏng?	
	\choice
	{Khoảng cách giữa các phân tử rất lớn so với kích thước của chúng}
	{\True Lực tương tác phân tử yếu hơn lực tương tác phân tử ở thể rắn}
	{Không có thể tích và hình dạng riêng xác định}
	{Các phân tử dao động xung quanh vị trí cân bằng xác định}
	\loigiai{
		\begin{itemize}
			\item Khoảng chất lỏng, các phân tử ở xa nhau hơn so với các phân	tử trong chất rắn nên lực tương tác giữa các phân tử chất lỏng nhỏ hơn trong chất
			rắn
		\end{itemize}
		
	}
	\haidong{3}\end{ex}

\begin{ex}
	Trong chuyển động nhiệt, các phân tử của chất lỏng
	\choice
	{chuyển động hoàn toàn hỗn loạn}
	{chuyển động hỗn loạn quanh vị trí cân bằng xác định}
	{dao động xung quanh các vị trí cân bằng cố định}
	{\True dao động quanh vị trí cân bằng nhưng những vị trí cân bằng này không cố định mà di chuyển}
	\loigiai{
		\begin{itemize}
			\item Các phân tử của chất lỏng dao động xung quanh các vị trí cân bằng và các vị trí cân bằng này lại có thể
			thay đổi.
		\end{itemize}
	}
	\haidong{2}\end{ex}

\begin{ex}
	Đặc điểm nào sau đây là đặc điểm của chất lỏng?
	\choice
	{Có hình dạng của phần bình chứa nó, có thể rót được và chảy tràn trên bề mặt, dễ nén}
	{\True Có hình dạng của phần bình chứa nó, có thể rót được và chảy tràn trên bề mặt, rất khó nén}
	{Có hình dạng cố định, có thể rót được và chảy tràn trên bề mặt, rất khó nén}
	{Có hình dạng của phần bình chứa nó, dễ dàng lan tỏa trong không gian, dễ nén}
	\loigiai{
		\begin{itemize}
			\item Do lực liên kết giữa các phân tử yếu hơn trong chất rắn nên các chất ở thể lỏng có thể tích xác định nhưng không có hình dạng riêng mà có hình dạng của phần bình chứa nó.
		\end{itemize}
	}
	\haidong{2}\end{ex}

\begin{ex}
	Chất lỏng có thể tích xác định, nhưng hình dạng không xác định là do trong chất lỏng có những đặc điểm nào sau đây?
	\choice
	{Lực liên kết giữa các phân tử chất lỏng là rất lớn, các phân tử chỉ dao động không ngừng quanh một vị trí xác định}
	{Lực liên kết giữa các phân tử chất lỏng là rất yếu, các phân tử dao động tự do về mọi phía}
	{\True Lực liên kết giữa các phân tử trong chất lỏng yếu hơn so với chất rắn, các phân tử khi dao động còn có thể trượt lên nhau}
	{Lực liên kết giữa các phân tử chất lỏng là rất yếu, các phân tử dao động xung quanh một vị trí xác định}
	\loigiai{
		\begin{itemize}
			\item Lực liên kết giữa các phân tử yếu hơn trong chất rắn nên các chất ở thể lỏng có thể tích xác định nhưng không có hình dạng riêng mà có hình dạng
			của phần bình chứa nó
		\end{itemize}
	}
	\haidong{1}\end{ex}



\begin{ex} \nguon{4}
	Chất lỏng \textbf{không} có hình dạng xác định vì các phân tử chất lỏng
	\choice
	{dao động tại các vị trí cân bằng xác định}
	{có thể chuyển động phân tán ra xa nhau}
	{\True dao động quanh các vị trí cân bằng có thể dịch chuyển được}
	{có thể chuyển động tự do}
	\loigiai{
		Chất lỏng có tính linh động vì các phân tử của nó dao động quanh các vị trí cân bằng có thể di chuyển được, giúp chất lỏng dễ dàng thay đổi hình dạng theo hình dạng của vật chứa.
	}
\end{ex}

\begin{ex} \nguon{4}
	Lực tương tác giữa các phân tử trong chất lỏng
	\choice
	{rất lớn và giữ cho các phân tử chất lỏng dao động tại vị trí cân bằng xác định}
	{rất yếu nên chất lỏng không có thể tích xác định}
	{\True vẫn đủ để giữ các phân tử không chuyển động phân tán ra xa nhau}
	{rất yếu nên chất lỏng rất dễ bị nén}\loigiai{
		Lực tương tác giữa các phân tử trong chất lỏng không mạnh như trong chất rắn nhưng đủ lớn để giữ cho các phân tử không chuyển động phân tán và giữ cho chất lỏng có thể tích xác định.
	}
\end{ex}

\begin{ex} \nguon{4}
	Chất lỏng có thể tích xác định vì
	\choice
	{các phân tử cấu tạo nên chất lỏng có thể chuyển động hỗn độn, tự do}
	{lực liên kết giữa các phân tử trong chất lỏng là rất yếu}
	{\True lực tương tác giữa các phân tử vẫn đủ để giữ các phân tử không chuyển động phân tán ra xa nhau}
	{vị trí cân bằng của các phân tử có thể chuyển động từ nơi này đến nơi khác}\loigiai{
		Lực tương tác giữa các phân tử trong chất lỏng đủ mạnh để giữ chúng lại gần nhau, tạo ra thể tích xác định cho chất lỏng mặc dù hình dạng của nó có thể thay đổi.
	}
\end{ex}

\begin{ex} \nguon{4}
	Chuyển động nào sau đây là chuyển động của riêng các phân tử ở thể lỏng?
	\choice
	{Chuyển động hỗn loạn không ngừng}
	{Dao động xung quanh các vị trí cân bằng cố định}
	{Chuyển động hoàn toàn tự do}
	{\True Dao động xung quanh các vị trí cân bằng không cố định}\loigiai{
		Các phân tử trong chất lỏng dao động quanh các vị trí cân bằng không cố định, cho phép chúng di chuyển linh hoạt và dễ dàng thay đổi vị trí.
	}
\end{ex}
\begin{ex} \nguon{1}
	Khi đổ 50 lít rượu vào 50 lít nước thì thu được bao nhiêu lít hỗn hợp?
	\choice
	{100 lít}
	{50 lít}
	{Lớn hơn 100 lít}
	{\True Nhỏ hơn 100 lít}
	\loigiai{
		\begin{itemize}
			\item Khi đổ 50 lít rượu vào 50 lít nước thì thu được \textbf{nhỏ hơn 100 lít} hỗn hợp. 
			\item Nguyên nhân là do các phân tử rượu và nước có thể xen kẽ và chen vào các khoảng trống giữa chúng, làm thể tích tổng giảm so với khi đo riêng rẽ từng phần. 
			\item Đây là một hiện tượng phổ biến khi trộn các chất có liên kết phân tử khác nhau.\end{itemize}
	}
\end{ex}
\begin{ex}
	immini[thm]{Hình bên là hình ảnh phóng to của bề mặt bàn. Hãy sử dụng mô hình động học phân tử để giải thích vì sao chất lỏng di chuyển trên mặt bàn dễ dàng hơn so với chất rắn}{	includegraphics[width=8cm]{images/caibannew}}
	\loigiai{
		\begin{itemize}
			\item Chất lỏng có thể di chuyển trên mặt bàn dễ dàng hơn so với chất rắn vì trong chất lỏng, các phân tử có sự tự do chuyển động tương đối, cho phép chúng chảy và thay đổi hình dạng. Ngược lại, các phân tử trong chất rắn được sắp xếp cứng nhắc và cố định tại các vị trí nhất định, do đó chất rắn không thể di chuyển hay chảy dễ dàng như chất lỏng.
		\end{itemize}
	}
\end{ex}

\loai{Loại 12: Cấu trúc chất khí}


\begin{ex} %[1P5Y1-1][Loai1]
	Phát biểu nào dưới đây là nguyên nhân gây ra chuyển động hoàn toàn hỗn loạn của các phân tử chất khí?
	\choice
	{Các phân tử luôn va chạm với nhau}
	{\True Lực tương tác giữa các phân tử rất yếu}
	{Khoảng cách giữa các phân tử nhỏ}
	{Khoảng cách giữa các phân tử không giống nhau}
	\loigiai{
		\begin{itemize}
			\item Ở thể khí, lực tương tác giữa các phân tử rất yếu nên các phân tử chuyển động hoàn toàn hỗn loạn.
		\end{itemize}
	}%<MyLT2>
	\haidong{2}\end{ex}

\begin{ex}%[1P5Y1-1][Loai2]
	Tính chất nào sau đây \textbf{không} phải là của phân tử vật chất ở thể khí?
	\choice
	{Chuyển động hỗn loạn}
	{Chuyển động không ngừng}
	{Chuyển động tự do về mọi phía}
	{\True Chuyển động hỗn loạn xung quanh các vị trí cân bằng cố định}
	\loigiai{
		\begin{itemize}
			\item Ở trạng thái khí, các phân tử chuyển động hỗn loạn và không có vị trí cân bằng.
		\end{itemize}
	} 
	%<MyLT3>
	\haidong{2}\end{ex}



\begin{ex}%[1P5Y1-1][Loai3]
	Phát biểu nào sau đây là đúng về chất khí?
	\choice
	{\True Chất khí không có hình dạng và thể tích riêng}
	{Chất khí luôn chiếm một phần thể tích của bình chứa}
	{Rất khó để nén được chất khí}
	{Chất khí có thể tích xác định}
	\loigiai{
		\begin{itemize}
			\item Do lực tương tác giữa các phân tử chất khí quá yếu nên chất khí không có hình dạng và thể tích riêng.
		\end{itemize}
	} 
	%<MyLT2>
	\haidong{3}\end{ex}

%\begin{ex} %[1P5Y1-1][Loai3]
%	Các tính chất nào sau đây là của phân tử chất khí?	
%	\choice
%	{Có thể tích xác định}
%	{\True Chuyển động càng nhanh khi nhiệt độ càng cao}
%	{Dao động quanh vị trí cân bằng}
%	{Luôn luôn tương tác với các phân tử khác}
%	\loigiai{
%		\begin{itemize}
%			\item Các phân tử chuyển động hỗn loạn không ngừng, nhiệt độ càng cao thì vận tốc chuyển động hỗn loạn càng lớn.
%		\end{itemize}
%	}%<MyLT1>
%	\haidong{3}\end{ex}

\begin{ex}%[1P5Y1-1][Loai3]
	Phát biểu nào sau đây là \textbf{sai} khi nói về chất khí?
	\choice
	{Lực tương tác giữa các phân tử là rất yếu}
	{\True Các phân tử khí ở rất gần nhau}
	{Chất khí không có hình dạng và thể tích riêng}
	{Chất khí luôn chiếm toàn bộ thể tích bình chứa và nén được dễ dàng}
	\loigiai{
		\begin{itemize}
			\item Phát biểu “Các phân tử khí ở rất gần nhau” là sai. Thực ra, các phân tử khí ở rất xa nhau (khoảng cách giữa các phân tử lớn gấp hàng chục lần kích thước của chúng).
		\end{itemize}
	}%<MyLT1>
	\haidong{3}\end{ex}

\begin{ex}
	Các phân tử khí chuyển động hỗn loạn không ngừng vì
	\choice
	{phân tử khí không có khối lượng}
	{khoảng cách giữa các phân tử khí quá gần nhau}
	{\True lực tương tác giữa các phân tử quá nhỏ}
	{các phân tử khí luôn đẩy nhau}
	\loigiai{
		\begin{itemize}
			\item Khoảng cách giữa các phân tử rất lớn so với kích
			thước của chúng. Lực tương tác giữa các phân tử hầu như không đáng
			kể (trừ khi va chạm nhau) nên chúng chuyển động hỗn loạn không ngừng về mọi phía
		\end{itemize}
	}
	\haidong{2}\end{ex}

\begin{ex}
	Tính chất nào sau đây \textbf{không} phải là tính chất của chất ở thể khí?
	
	\choice
	{\True Có hình dạng và thể tích riêng}
	{Các phân tử chuyển động hoàn toàn hỗn độn}
	{Có thể nén được dễ dàng}
	{Có lực tương tác phân tử nhỏ hơn lực tương tác phân tử ở thể rắn và thể lỏng}
	\loigiai{
		\begin{itemize}
			\item Các chất ở thể khí không có thể tích và hình dạng riêng mà có thể tích và hình dạng của bình chứa.
		\end{itemize}
	}
	\haidong{3}\end{ex}

\begin{ex}
	Chất khí \textbf{không} có hình dạng và thể tích riêng là vì
	
	\choice
	{khoảng cách giữa các phân tử rất gần, lực tương tác giữa các phân tử chất khí rất mạnh}
	{khoảng cách giữa các phân tử rất gần, lực tương tác giữa các phân tử chất khí rất yếu}
	{khoảng cách giữa các phân tử rất xa, lực tương tác giữa các phân tử chất khí rất mạnh}
	{\True khoảng cách giữa các phân tử rất xa, lực tương tác giữa các phân tử chất khí rất yếu}
	\loigiai{
		\begin{itemize}
			\item  Khoảng cách giữa các phân tử rất lớn so với kích
			thước của chúng nên lực tương tác giữa các phân tử hầu như không đáng
			kể (trừ khi va chạm nhau) dẫn tới hệ quả trên.
		\end{itemize}
	}
	\haidong{4}\end{ex}

\begin{ex}
	Chất khí luôn luôn chiếm hết thể tích của bình chứa bởi vì
	
	\choice
	{\True lực liên kết giữa các phân tử chất khí rất yếu nên chúng chuyển động tự do về mọi phía}
	{các phân tử chất khí luôn chuyển động nhiệt không ngừng}
	{trong chất khí có quá nhiều phân tử}
	{các phân tử chất khí luôn luôn đẩy nhau ra xa nên chúng cách nhau càng xa càng tốt}
	\loigiai{
		\begin{itemize}
			\item  Khoảng cách giữa các phân tử rất lớn so với kích
			thước của chúng nên lực tương tác giữa các phân tử hầu như không đáng
			kể (trừ khi va chạm nhau) nên chúng chuyển động tự do về mọi phía
		\end{itemize}
	}
	\haidong{3}\end{ex}

\begin{ex}
	immini[thm]{Hình bên mô tả cấu trúc phân tử ở thể nào dưới đây?
		\choice
		{Thể lỏng}
		{\True Thể khí}
		{Thể rắn}
		{Plasma}}
	{\begin{tikzpicture}[scale=.8,thick,>=stealth']
			\draw[fill=blue!10] (0,0) -- (0,-2) -- (5,-2) -- (5,0) -- (0,0);
			\shade[ball color=blue!70] (0.4,-0.4) circle (1.3ex);
			%\shade[ball color=blue!70] (1,-1.5) circle (1.3ex);
			\shade[ball color=blue!70] (2,-1.6) circle (1.3ex);
			\shade[ball color=blue!70] (3.2,-0.3) circle (1.3ex);
			%\shade[ball color=blue!70] (4.25,-0.75) circle (1.3ex);
			\shade[ball color=blue!70] (4.5,-1.5) circle (1.3ex);
			\shade[ball color=blue!70] (5.5,-0.25) circle (1.3ex);
			\shade[ball color=blue!30] (6,-1) circle (1.3ex);
			\shade[ball color=blue!30] (5.5,-1.75) circle (1.3ex);
			\draw[dashed,thin,->] (5.6,-0.4) -- (5.9,-0.8);
			\draw [dashed,thin,->] (6,-1) -- (5.6,-1.6);
			\draw [dashed,thin,->] (5.5,-1.8) -- (6.1,-2);
	\end{tikzpicture}}
	\loigiai{
		\begin{itemize}
			\item  Khoảng cách giữa các phân tử rất lớn nên đây là thể khí
		\end{itemize}
	}
	\haidong{2}\end{ex}

\vspace*{0.2 cm}



\begin{ex} \nguon{3}
	Một bình nước hoa được mở nắp ở cuối lớp học nhưng ở mọi vị trí của lớp hầu như bạn nào cũng ngửi thấy được. Điều này thể hiện tính chất nào của thể khí?
	\choice
	{Dễ dàng nén được}
	{Không có hình dạng xác định}
	{\True Có thể lan toả trong không gian theo mọi hướng}
	{Không chảy được}
	\loigiai{
		Tính chất "Có thể lan toả trong không gian theo mọi hướng" của chất khí giải thích vì sao mọi người trong lớp đều ngửi thấy được mùi nước hoa khi bình nước hoa được mở nắp. Điều này là do các phân tử khí của nước hoa phân tán và di chuyển khắp phòng.
	}
\end{ex}


\begin{ex} \nguon{4}
	Phát biểu nào sau đây về các phân tử của chất khí là đúng?
	\choice
	{Các phân tử chất khí liên kết chặt chẽ với nhau}
	{Các phân tử chất khí dao động quanh vị trí cân bằng cố định}
	{\True Các phân tử chất khí chuyển động hỗn loạn, không ngừng về mọi phía}
	{Các phân tử chất khí được sắp xếp rất gần nhau}
	\loigiai{
		
		Các phân tử chất khí luôn chuyển động hỗn loạn và không ngừng về mọi phía, điều này làm cho chất khí có thể lan tỏa đều trong không gian chứa nó.
	}
\end{ex}

\begin{ex} \nguon{4}
	Tính chất nào sau đây \textbf{không} phải là của phân tử của các chất ở thể khí?
	\choice
	{Các phân tử liên kết với nhau rất yếu}
	{Khoảng cách trung bình của các phân tử xa nhau}
	{Chuyển động hỗn loạn và không ngừng}
	{\True Chuyển động hỗn loạn xung quanh các vị trí cân bằng cố định}
	\loigiai{
		
		Các phân tử khí không dao động xung quanh các vị trí cân bằng cố định. Thay vào đó, chúng chuyển động hỗn loạn và không ngừng trong không gian.
	}
\end{ex}

\begin{ex} \nguon{4}
	Chất khí \textbf{không} có thể tích riêng là vì
	\choice
	{các phân tử chất khí dao động quanh các vị trí cân bằng có thể di chuyển được}
	{\True các phân tử chất khí chuyển động hỗn loạn về mọi phía}
	{lực đẩy giữa các phân tử chất khí rất mạnh}
	{khoảng cách giữa các phân tử chất khí rất nhỏ so với kích thước của chúng}
	\loigiai{
		
		Các phân tử chất khí chuyển động hỗn loạn theo mọi hướng, khiến cho chất khí không có thể tích cố định hay hình dạng cố định.
	}
\end{ex}


\loai{Loại 13: Tổng hợp các thể vật chất đơn giản}
\begin{ex}%[1P5B1-1][Loai2]
	Phát biểu nào sau đây là \textbf{sai}?
	\choice
	{Các chất được cấu tạo một cách gián đoạn}
	{Lực tương tác giữa các phân tử ở thể rắn lớn hơn lực tương tác giữa các phân tử ở thể lỏng và thể khí}
	{\True Các phân tử chất rắn dao động quanh vị trí cân bằng không cố định}
	{Giữa các phân tử đồng thời có cả lực hút và lực đẩy}
	\loigiai{
		\begin{itemize}
			\item Ở trạng thái rắn, các nguyên tử, phân tử dao động quanh các vị trí cân bằng cố định.
		\end{itemize}
	} %<MyLT1>
	\haidong{2}\end{ex}

\begin{ex}%[1P5Y1-1][Loai1]
	Chất nào rất khó nén được?
	\choice
	{Chỉ có chất rắn}
	{\True Chất rắn, chất lỏng}
	{Chất khí, chất lỏng}
	{Chất khí, chất rắn}
	\loigiai{
		\begin{itemize}
			\item Do khoảng cách các phân tử ở gần nhau nên chất lỏng và chất rắn gần như không nén được.
		\end{itemize}
	} 
	\haidong{2}\end{ex}


\begin{ex} \nguon{1}
	Ở nhiệt độ phòng, chất nào ở thể lỏng?
	\choice
	{Chì}
	{Kẽm}
	{\True Thủy ngân}
	{Thiếc}
	\loigiai{
		\begin{itemize}
			\item Chất ở thể lỏng ở nhiệt độ phòng là: \emph{Thủy ngân}.
			
			\item Giải thích chi tiết:
			\begin{itemize}
				\item \textbf{Chì}: Ở nhiệt độ phòng, chì là chất rắn.
				\item \textbf{Kẽm}: Ở nhiệt độ phòng, kẽm là chất rắn.
				\item \textbf{Thủy ngân}: Đúng. Thủy ngân là một chất đặc biệt, ở thể lỏng tại nhiệt độ phòng (khoảng 25°C).
				\item \textbf{Thiếc}: Ở nhiệt độ phòng, thiếc là chất rắn.
		\end{itemize}\end{itemize}
	}
\end{ex}




\loai{Loại 14: So sánh đặc điểm của các chất}
%\begin{ex}%[1P5Y1-1][Loai2]
%	Trong các tính chất sau, tính chất nào \textbf{không} chỉ là của chất khí?
%	\choice
%	{\True Không có hình dạng cố định}
%	{Chiếm toàn bộ thể tích của bình chứa}
%	{Tác dụng lực lên mọi phần diện tích bình chứa}
%	{Thể tích giảm đáng kể khi tăng áp suất}
%	\loigiai{
%		\begin{itemize}
%			\item Chất lỏng cũng không có hình dạng cố định.
%		\end{itemize}
%	} 
%	\haidong{3}\end{ex}



%\begin{ex}[KNTT] Hãy giải thích các đặc điểm sau đây của thể khí, thể rắn, thể lỏng:
%	\begin{itemize}
%		\item Chất khí không có hình dạng và thể tích riêng, luôn chiếm toàn bộ thể tích bình chứa và có thể nén được dễ dàng.
%		\item Vật ở thể rắn có thể tích và hình dạng riêng, rất khó nén.
%		\item Vật ở thể lỏng có thể tích riêng nhưng không có hình dạng riêng.
%	\end{itemize}
%	\loigiai{
%		\begin{itemize}
%			\item Ở thể khí, các phân tử ở xa nhau (khoảng cách giữa các phân tử lớn gấp hàng chục lần kích thước của chúng). Lực tương tác giữa các phân tử rất yếu nên các phân tử chuyển động hoàn toàn hỗn loạn. Do đó, chất khí không có hình dạng và thể tích riêng. Chất khí luôn chiếm toàn bộ thể tích của bình chứa và có thể nén được dễ dàng.
%			\item Ở thể rắn, các phân tử ở gần nhau (khoảng cách giữa các phân tử chỉ vào cỡ kích thước của chúng). Lực tương tác giữa các phân tử chất rắn rất mạnh nên giữ được các phân tử này ở các vị trí xác định và làm cho chúng chỉ có thể dao động xung quanh các vị trí cân bằng xác định này. Do đó, các vật rắn có thể tích và hình dạng riêng xác định, rất khó nén.
%			\item Lực tương tác giữa các phân tử ở thể lỏng lớn hơn lực tương tác giữa các phân tử ở thể khí nên giữ được các phân tử không chuyển động phân tán ra xa nhau. Nhờ đó, chất lỏng có thể tích riêng xác định. Tuy nhiên, lực này chưa đủ lớn như trong chất rắn để giữ các phân tử ở những vị trí xác định. Các phân tử ở thể lỏng cũng dao động xung quanh các vị trí cân bằng, nhưng những vị trí này không cố định mà di chuyển. Do đó chất lỏng không có hình dạng riêng mà có hình dạng của phần bình chứa nó.
%		\end{itemize}
%		
%	}
%	\haidong{16}\end{ex}


\begin{ex} \nguon{1}
	Kim cương và than chì có tính chất vật lí khác nhau là do chúng khác nhau về
	\choice
	{\True cấu trúc mạng tinh thể}
	{bản chất các hạt tạo thành tinh thể}
	{loại liên kết giữa các hạt trong tinh thể}
	{kích thước tinh thể}
	\loigiai{
		\begin{itemize}
			\item Kim cương và than chì đều được tạo thành từ các nguyên tử cacbon, nhưng chúng có cấu trúc tinh thể khác nhau. 
			\item Trong kim cương, mỗi nguyên tử carbon liên kết với bốn nguyên tử carbon khác xung quanh theo cấu trúc tứ diện, tạo nên một mạng lưới ba chiều rất cứng và bền. 
			\item Trong khi đó, trong than chì, các nguyên tử carbon liên kết tạo thành tấm phẳng với các liên kết mạnh trong từng lớp, nhưng các lớp này lại liên kết với nhau bằng các lực yếu hơn (lực Van der Waals), cho phép chúng trượt lên nhau dễ dàng. 
			\item Sự khác biệt trong cấu trúc tinh thể này dẫn đến các tính chất vật lí rất khác nhau giữa kim cương và than chì.\end{itemize}
	}
\end{ex}


\begin{ex} \nguon{1}
	Phát biểu nào sau đây \textbf{không} đúng?
	\choice
	{Các phân tử chất rắn dao động xung quanh các vị trí cân bằng xác định}
	{Các phân tử chất lỏng dao động xung quanh các vị trí cân bằng có thể di chuyển được}
	{Các phân tử chất khí không dao động xung quanh các vị trí cân bằng}
	{\True Các phân tử chất rắn, chất lỏng và chất khí đều chuyển động hỗn độn giống nhau}
	\loigiai{
		\begin{itemize}
			\item Phát biểu \textbf{'Các phân tử chất rắn, chất lỏng và chất khí đều chuyển động hỗn độn như nhau'} là không đúng. 
			\item Trong chất rắn, các phân tử dao động xung quanh các vị trí cân bằng cố định do cấu trúc mạng tinh thể vững chắc.
			\item  Trong chất lỏng, các phân tử dao động xung quanh các vị trí cân bằng nhưng có thể di chuyển tương đối tự do so với nhau.
			\item  Trong chất khí, các phân tử chuyển động hỗn độn theo mọi hướng và không dao động xung quanh vị trí cân bằng cố định mà chuyển động tự do với khoảng cách lớn giữa các phân tử.
			\item  Do đó, mức độ và tính chất của chuyển động ở các trạng thái này là khác nhau.\end{itemize}
	}
\end{ex}

\begin{ex} \nguon{4}
	Khi nói về khoảng cách trung bình giữa các phân tử trong chất rắn, chất lỏng, chất khí. Phát biểu nào sau đây là đúng?
	\choice
	{Khoảng cách giữa các phân tử trong chất lỏng xa hơn so với các phân tử trong chất khí}
	{Khoảng cách giữa các phân tử trong chất rắn xa hơn so với các phân tử trong chất lỏng}
	{\True Khoảng cách giữa các phân tử trong chất lỏng gần hơn so với các phân tử trong chất khí}
	{Khoảng cách giữa các phân tử trong chất khí gần hơn so với các phân tử trong chất lỏng}
	\loigiai{
		\begin{itemize}
			\item Trong chất lỏng, các phân tử ở gần nhau hơn so với trong chất khí nhưng xa nhau hơn so với trong chất rắn.
			\item Chất khí có khoảng cách giữa các phân tử lớn nhất, tạo nên tính chất dễ nén.
			\item Chất rắn có khoảng cách giữa các phân tử nhỏ nhất, tạo nên cấu trúc vững chắc.
		\end{itemize}
	}
\end{ex}

\begin{ex} \nguon{4}
	Gọi lực liên kết giữa các phân tử trong chất rắn, chất lỏng, chất khí lần lượt là $F_1, \ F_2, \ F_3$ thì
	\choice
	{$F_1=F_2=F_3$}
	{\True $F_1>F_2>F_3$}
	{$F_1<F_2=F_3$}
	{$F_1>F_2=F_3$}
	\loigiai{
		\begin{itemize}
			\item Lực liên kết giữa các phân tử trong chất rắn ($F_1$) lớn hơn trong chất lỏng ($F_2$) và trong chất khí ($F_3$).
			\item Điều này là do phân tử trong chất rắn gắn kết rất chặt chẽ, trong khi ở chất lỏng và khí lực gắn kết yếu hơn.
		\end{itemize}
	}
\end{ex}
\begin{ex} \nguon{1}
	Chất lỏng được dẫn bằng các đường ống. Điều này là do yếu tố nào sau đây?
	\choice
	{\True Tính chất thay đổi hình dạng theo bình chứa}
	{Chất lỏng có cấu trúc mạng tinh thể}
	{Các phân tử chất lỏng chuyển động hỗn loạn} 
	{Các phân tử chất lỏng va chạm với thành bình}
	\loigiai{
		Thể lỏng có hình dạng thay đổi được theo hình dạng của phần bình chứa nó.
	}
\end{ex}
\begin{ex} \nguon{1}
	Phát biểu nào sau đây \textbf{không} đúng?
	\choice
	{Các phân tử chất lỏng có thể trượt lên nhau}
	{Các phân tử chất lỏng dao động xung quanh các vị trí cân bằng có thể di chuyển được}
	{Lực liên kết giữa các phân tử trong chất lỏng nhỏ hơn so với trong chất rắn}
	{\True Chất lỏng dễ dàng nén được còn chất rắn thì không}
	\loigiai{
		\begin{itemize}
			\item Phát biểu \textbf{'Chất lỏng dễ dàng nén được còn chất rắn thì không'} là không đúng. 
			\item  Các chất lỏng nói chung rất khó nén vì các phân tử trong chất lỏng đã xếp sát nhau, tương tự như trong chất rắn.
			\item  Trái lại, chất khí dễ dàng nén được vì khoảng cách giữa các phân tử rất lớn.
			\item  Các phát biểu khác đều đúng vì các phân tử chất lỏng có thể trượt lên nhau (cho phép chất lỏng chảy), dao động xung quanh các vị trí cân bằng có thể di chuyển được và lực liên kết giữa các phân tử chất lỏng nhỏ hơn so với chất rắn.\end{itemize}
	}
\end{ex}



\begin{ex} \nguon{4}
	Chất khí dễ bị nén hơn so với chất rắn và chất lỏng vì
	\choice
	{lực tương tác giữa các phân tử trong chất khí lớn hơn so với lực tương tác giữa các phân tử trong chất rắn và chất lỏng}
	{\True khoảng cách giữa các phân tử trong chất khí lớn hơn so với khoảng cách giữa các phân tử trong chất rắn và chất lỏng}
	{các phân tử trong chất khí ít chuyển động hơn so với các phân tử trong chất rắn và chất lỏng}
	{các phân tử trong chất khí có kích thước nhỏ hơn so với các phân tử trong chất rắn và chất lỏng}
	\loigiai{
		\begin{itemize}
			\item Khoảng cách giữa các phân tử trong chất khí lớn hơn so với khoảng cách giữa các phân tử trong chất rắn và chất lỏng. 
			\item Điều này có nghĩa là giữa các phân tử chất khí có nhiều không gian trống hơn, do đó khi nén, các phân tử có thể di chuyển gần nhau hơn mà không gặp phải lực cản lớn như trong chất rắn và chất lỏng.
			\item Ở chất rắn và chất lỏng, các phân tử đã nằm sát nhau hơn, nên rất khó để nén thêm nữa.
		\end{itemize}
	}
\end{ex}







\begin{ex}
	Tại sao quả bóng bay dù được buộc chặt nhưng để lâu ngày vẫn bị xẹp?	
	\choice
	{Vì khi thổi, không khí từ miệng vào bóng còn nóng, sau đó lạnh dần nên co lại}
	{Vì cao su là chất đàn hồi nên sau khi bị thổi căng nó tự động co lại}
	{Vì không khí nhẹ nên có thể chui qua chỗ buộc ra ngoài}
	{\True Vì giữa các phân tử của chất làm vỏ bóng có khoảng cách nên các phân tử không khí có thể qua đó thoát ra ngoài}
	\loigiai{
		\begin{itemize}
			\item Khi bơm quả bóng phồng lên. Vỏ quả bóng dãn ra làm khoảng cách giữa các phân tử của vỏ bóng tăng lên
			\item Lực liên kết giữa các phân tử khí bên trong bóng lỏng lẻo nên các phân tử khí này dễ dàng chui qua khoảng cách giữa các phân tử làm vỏ bóng ra ngoài
		\end{itemize}
		
	}
	\haidong{2}\end{ex}
	























\loai{Loại 16: Câu hỏi đúng sai}
\begin{ex} \nguon{1}
	Chuyển động của phân tử có đặc điểm nào sau đây?
	\choiceTF[t]
{Hoàn toàn do lực tương tác phân tử gây ra}
{Chuyển động không ngừng}
{Chuyển động càng nhanh thì nhiệt độ càng cao}
{Luôn dao động quanh vị trí cân bằng}
	\loigiai{
		\begin{itemize}
			\item Sai. Chuyển động của các phân tử không phải do lực tương tác phân tử gây ra mà do năng lượng nhiệt. Lực tương tác giữa các phân tử ảnh hưởng đến hướng và tốc độ của chuyển động, nhưng bản thân chuyển động là do năng lượng nhiệt.
			
			\item Đúng. Các phân tử luôn chuyển động không ngừng do chúng có năng lượng động học. Sự chuyển động này có thể là tịnh tiến, dao động hoặc quay tùy thuộc vào trạng thái của vật chất.
			
			\item Đúng. Chuyển động của các phân tử càng nhanh thì nhiệt độ của chất càng cao. Điều này là do nhiệt độ là thước đo trung bình của động năng phân tử.
			
			\item Sai. Các phân tử không luôn dao động quanh vị trí cân bằng. Điều này đúng chủ yếu đối với chất rắn, nơi các phân tử dao động quanh vị trí cố định của chúng. Trong chất khí các phân tử chuyển động tự do trong không gian.
		\end{itemize}
	}
\end{ex}
\begin{ex}
	Theo mô hình động học phân tử thì
	\choiceTF[t]
{giữa các phân tử không có khoảng cách}
{lực tương tác giữa các phân tử của vật ở thể rắn lớn hơn lực tương tác giữa các phân tử của vật ở thể lỏng, thể khí}
{các phân tử chất rắn dao động xung quanh các vị trí cân bằng không cố định}
{các phân tử chất lỏng dao động xung quanh các vị trí cân bằng không cố định}
	\loigiai{
		\begin{itemize}
			\item 	SAI. Giữa các phân tử có khoảng cách
			\item 	ĐÚNG
			\item 	SAI. Các phân tử chất rắn dao động xung quanh các vị trí cân bằng cố định.
			\item 	ĐÚNG
		\end{itemize}
	}
	\haidong{10}\end{ex}

\begin{ex}
	Theo mô hình động học phân tử thì
	\choiceTF[t]
{giữa các phân tử có lực tương tác. Độ lớn của những lực này phụ thuộc vào khoảng cách giữa chúng}
{các phân tử chuyển động không ngừng, chuyển động này càng nhanh thì nhiệt độ của vật càng cao}
{các chất được cấu tạo từ những hạt có kích thước rất nhỏ gọi là phân tử. Giữa các phân tử có khoảng cách}
{khi khoảng cách giữa các phân tử lớn thì lực đẩy mạnh hơn lực hút}
	
	\loigiai{
		\begin{itemize}
			\item ĐÚNG
			\item ĐÚNG
			\item ĐÚNG
			\item SAI. Khi khoảng cách giữa các phân tử lớn thì lực hút mạnh hơn lực đẩy.
		\end{itemize}
	}
	\haidong{10}\end{ex}

\begin{ex} \nguon{1}
	Theo mô hình động học phân tử về cấu tạo chất thì
	\choiceTF[t]
{các chất được cấu tạo từ các hạt riêng biệt là các phân tử}
{các phân tử đứng sát nhau và giữa chúng không có khoảng cách}
{các phân tử có lúc đứng yên, có lúc chuyển động}
{giữa các phân tử đồng thời có cả lực hút và lực đẩy}
	\loigiai{
		\begin{itemize}
			\item Đúng. Theo mô hình động học phân tử, mọi chất đều được cấu tạo từ các hạt riêng biệt như nguyên tử hay phân tử. Các hạt này là đơn vị cơ bản cấu tạo nên vật chất.
			
			\item Sai. Các phân tử không đứng sát nhau mà giữa chúng có khoảng cách. Khoảng cách này khác nhau tùy thuộc vào trạng thái của chất; ví dụ, trong chất khí, khoảng cách giữa các phân tử lớn hơn nhiều so với trong chất rắn và chất lỏng.
			
			\item Sai. Theo lý thuyết động học phân tử, các phân tử luôn chuyển động. Sự chuyển động của các phân tử có thể là dao động, chuyển động tịnh tiến hoặc chuyển động quay, nhưng chúng không bao giờ đứng yên hoàn toàn, kể cả ở nhiệt độ rất thấp.
			
			\item Đúng. Giữa các phân tử luôn tồn tại các lực tương tác lẫn nhau, bao gồm cả lực hút và lực đẩy. Lực hút giữ các phân tử lại với nhau, trong khi lực đẩy ngăn cản các phân tử chạm vào nhau quá gần. Những lực này đóng vai trò quan trọng trong việc xác định các tính chất vật lí của chất ở các trạng thái khác nhau.
		\end{itemize}
	}
\end{ex}
\begin{ex} \nguon{1}
Theo thuyết động học phân tử về cấu tạo chất thì các phân tử chất rắn 
	\choiceTF[t]
{có hình dạng cố định}
{chiếm toàn bộ thể tích của bình chứa}
{có lực tương tác phân tử lớn}
{các phân tử luôn dao động quanh vị trí cân bằng không cố định}
	\loigiai{
		\begin{itemize}
			\item Đúng. Chất rắn có hình dạng cố định.
			
			\item Sai. Chất rắn không chiếm toàn bộ thể tích của bình chứa mà giữ nguyên thể tích cố định của mình.
			
			\item Đúng. Lực tương tác giữa các phân tử trong chất rắn rất lớn, giữ cho các phân tử ở vị trí cố định.
			
			\item Sai. Các phân tử trong chất rắn dao động quanh các vị trí cân bằng cố định, không phải các vị trí không cố định.
		\end{itemize}
	}
\end{ex}
\begin{ex} \nguon{1}
	Theo thuyết động học phân tử về cấu tạo chất thì các phân tử của chất lỏng 
	\choiceTF[t]
{có khoảng cách giữa các phân tử nhỏ hơn nhiều so với trong chất khí}
{có lực tương tác giữa các phân tử khá lớn}
{không có hình dạng và thể tích cố định}
{các phân tử luôn dao động quanh vị trí cân bằng cố định}
	\loigiai{
		\begin{itemize}
			\item Đúng. Trong chất lỏng, khoảng cách giữa các phân tử nhỏ hơn nhiều so với trong chất khí, nhưng vẫn lớn hơn trong chất rắn.
			
			\item Đúng. Lực tương tác giữa các phân tử trong chất lỏng khá lớn, mạnh hơn so với trong chất khí nhưng yếu hơn so với trong chất rắn.
			
			\item Sai. Chất lỏng không có hình dạng cố định (nó thay đổi theo hình dạng của bình chứa) nhưng có thể tích cố định.
			
			\item Sai. Chỉ trong chất rắn, các phân tử mới dao động quanh vị trí cân bằng cố định. Trong chất lỏng, các phân tử có thể chuyển động tự do hơn, không chỉ dao động quanh vị trí cố định.
		\end{itemize}
	}
\end{ex}

\begin{ex} \nguon{1}
Theo thuyết động học phân tử về cấu tạo chất thì các phân tử của chất khí
	\choiceTF[t]
{có lực tương tác giữa các phân tử rất yếu}
{ở rất gần nhau}
{có hình dạng và thể tích riêng}
{luôn luôn chiếm toàn bộ thể tích bình chứa và có thể nén được dễ dàng}
	\loigiai{
		\begin{itemize}
			\item Đúng. Trong chất khí, lực tương tác giữa các phân tử rất yếu, cho phép chúng chuyển động tự do và lan tỏa ra khắp không gian.
			
			\item Sai. Các phân tử khí ở xa nhau hơn so với các phân tử trong chất lỏng hoặc chất rắn, đó là lý do chất khí có mật độ nhỏ và có thể dễ dàng nén lại.
			
			\item Sai. Chất khí không có hình dạng và thể tích riêng. Chúng sẽ chiếm toàn bộ thể tích của bình chứa mà chúng được đặt vào.
			
			\item Đúng. Chất khí luôn luôn chiếm toàn bộ thể tích bình chứa và có thể nén được dễ dàng vì khoảng cách giữa các phân tử rất lớn, cho phép chúng bị nén lại dưới áp lực.
		\end{itemize}
	}
\end{ex}
\begin{ex}[CTST] 
	Mở lọ nước hoa và đặt ở một góc trong phòng, một lúc sau, người trong phòng có thể ngửi thấy mùi nước hoa. Hiện tượng này được giải thích do
\choiceTF[t]
{gió ở quạt thổi các phân tử nước hoa đến chỗ người}
{trong hỗn hợp nước hoa có chứa cồn}
{người hít thở nên các hạt nước hoa chuyển động lại gần}
{các phân tử nước hoa chuyển động hỗn loạn và khuếch tán ra khắp phòng}
	\loigiai{
		\begin{itemize}
			\item Sai b) Sai c) Sai d) Đúng
			\item Các phân tử khí nước hoa chuyển động hỗn loạn không ngừng trong không khí, chúng khuếch tán trong không khí một cách nhanh chóng, cho nên một lúc sau người trong phòng có thể ngửi thấy mùi nước hoa.
		\end{itemize}
		
	}
	\haidong{8}\end{ex}
\newpage 
\begin{ex} \nguon{4}
	Mô hình động học phân tử có thể dùng để giải thích cấu trúc của các chất rắn, chất lỏng và chất khí. Các chất được cấu tạo từ các nguyên tử, phân tử và chúng chuyển động không ngừng. Sự sắp xếp và tương tác khác nhau giữa các phân tử sẽ tạo nên các tính chất khác nhau cho chất rắn, chất lỏng, chất khí.
	\choiceTF[t]
{Các chất rắn có hình dạng và thể tích xác định vì các phân tử cấu tạo nên nó đứng yên tại chỗ}
{Các chất khí dễ nén hơn các chất lỏng vì khoảng cách trung bình giữa các phân tử chất khí lớn hơn khoảng cách trung bình giữa các phân tử chất lỏng}
{Chất lỏng không có hình dạng xác định vì các phân tử dao động quanh những vị trí cân bằng có thể di chuyển được}
{Tất cả các chất rắn đều có cấu trúc mạng tinh thể nên có hình dạng xác định}
	\loigiai{
		\begin{itemize}
			\item[a)] Sai. Các phân tử vẫn dao động quanh vị trí cân bằng cố định.
			\item[b)] Đúng.
			\item[c)] Đúng.
			\item[d)] Sai. Chỉ có chất rắn kết tinh mới có cấu trúc mạng tinh thể, chất rắn vô định hình (thủy tinh, nhựa, sôcôla...) không có cấu trúc mạng tinh thể.
	\end{itemize}}
\end{ex}

\begin{ex} \nguon{4}
	Theo mô hình động học phân tử thì các phân tử cấu tạo nên các vật có sự tương tác với nhau. Sự khác biệt về độ lớn của lực tương tác giữa các phân tử trong chất rắn, chất lỏng, chất khí dẫn đến sự khác nhau về cấu trúc của chúng.
	\choiceTF[t]
{Lực tương tác giữa các phân tử chất rắn mạnh hơn lực tương tác giữa các phân tử chất khí}
{Lực tương tác giữa các phân tử chất rắn rất mạnh nên giữ các phân tử chỉ dao động quanh các vị trí cân bằng cố định, vì vậy chất rắn có hình dạng xác định}
{Các phân tử chất lỏng tương tác với nhau rất yếu nên các phân tử chất lỏng có thể chuyển động tự do ra xa nhau, vì vậy chất lỏng không có hình dạng xác định}
{Các chất khí có khoảng cách giữa các phân tử rất lớn so với kích thước của chúng nên lực tương tác giữa các phân tử gần như không đáng kể}
	\loigiai{
		\begin{itemize}
			\item Đúng.
			\item Đúng
			\item Sai. Lực tương tác giữa các phân tử chất lỏng vẫn đủ để giữ các phân tử chất lỏng không chuyển động phân tán ra xa nhau.
			\item Đúng.
		\end{itemize}
	}
\end{ex}




\begin{ex}
	Xét cấu trúc của chất rắn thì
	\choiceTF[t]
{các phân tử chất rắn ở rất gần nhau và sắp xếp một cách chặt chẽ, có trật tự}
{chất rắn có thể tích và hình dạng không xác định}
{muối ăn và kim cương là chất rắn vô định hình}
{chất rắn kết tinh có cấu trúc tinh thể}
	\loigiai{
		\begin{itemize}
			\item Sai. Chỉ có chất rắn kết tinh thì các phân tử mới có trật tự.
			\item SAI. Chất rắn có thể tích và hình dạng xác định.
			\item SAI. Muối ăn và kim cương là chất rắn kết tinh
			\item ĐÚNG
		\end{itemize}
	}
	\haidong{6}\end{ex}

\begin{ex}
	Xét cấu trúc của chất lỏng thì
	\choiceTF[t]
{khoảng cách trung bình giữa các phân tử trong chất lỏng lớn hơn khoảng cách trung bình giữa các phân tử trong chất rắn và nhỏ hơn khoảng cách trung bình của các phân tử trong chất khí}
{các phân tử chất lỏng dao động xung quanh vị trí cân bằng cố định}
{chất lỏng có thể tích xác định nhưng không có hình dạng xác định mà hình dạng của nó phụ thuộc vào hình dạng của phần bình chứa nó}
{lực tương tác giữa các phân tử ở thể lỏng lớn hơn lực tương tác giữa các phân tử ở thể khí}
	\loigiai{
		\begin{itemize}
			\item ĐÚNG
			\item SAI. Các phân tử chất lỏng dao động xung quanh vị trí cân bằng không cố định.
			\item ĐÚNG
			\item ĐÚNG
		\end{itemize}
	}
	\haidong{8}\end{ex}

\begin{ex}
	Xét cấu trúc của chất khí thì
	\choiceTF[t]
{khoảng cách giữa các phân tử chất khí rất lớn so với kích thước của chúng}
{trừ khi va chạm, lực tương tác giữa các phân tử khí rất nhỏ, hầu như không đáng kể}
{một lượng không khí xác định luôn có thể tích và hình dạng riêng xác định}
{các phân tử chất khí chuyển động hỗn loạn, không ngừng về mọi phía, chiếm toàn bộ không gian bình chứa}
	\loigiai{
		\begin{itemize}
			\item ĐÚNG
			\item ĐÚNG
			\item SAI. Chất khí không có thể tích và hình dạng xác định
			\item ĐÚNG
		\end{itemize}
	}
	\haidong{5}\end{ex}



%\begin{ex} \nguon{1}
%	Khi tăng nhiệt độ của một chất thì
%	\choiceTF[t]
{tất cả các chất đều nở ra}
{các phân tử sẽ chuyển động nhanh hơn, hỗn loạn hơn, dao động quanh vị trí cân bằng được tăng cường}
{lực liên kết giữa các phân tử suy giảm}
{các chất sẽ chuyển thể}
%	\loigiai{
%		\begin{itemize}
%			\item Sai. Không phải tất cả các chất đều nở ra khi tăng nhiệt độ; ví dụ, nước từ 0°C đến 4°C sẽ co lại.
%			
%			\item Đúng. Khi nhiệt độ tăng, các phân tử sẽ chuyển động nhanh hơn và hỗn loạn hơn, làm tăng cường dao động quanh vị trí cân bằng.
%			
%			\item Đúng. Tăng nhiệt độ làm giảm lực liên kết giữa các phân tử, dẫn đến khả năng phá vỡ các liên kết này.
%			
%			\item Sai. Tăng nhiệt độ có thể dẫn đến chuyển thể chứ không chắc chắn làm cho tất cả các chất chuyển thể.
%		\end{itemize}
%	}
%\end{ex}

%\begin{ex} \nguon{3}
%	Cho biết các phát biểu sau đây đúng hay sai?
%	\choiceTF[t]
{Khoảng cách giữa các phân tử lớn thì lực liên kết giữa chúng yếu}
{Các phân tử sắp xếp càng có trật tự thì lực liên kết giữa chúng càng mạnh}
{Vật ở thể lỏng có thể tích và hình dạng riêng}
{Vật ở thể rắn có thể tích và hình dạng riêng, rất khó nén}
%	\loigiai{
%		\begin{itemize}
%			\item Đúng
%			\item Đúng
%			\item Sai
%			\item Đúng
%		\end{itemize}
%		Vật ở thể lỏng có thể tích xác định, nhưng không có hình dạng xác định.
%	}
%\end{ex}
%
%\begin{ex} \nguon{3}
%	Cho biết các phát biểu sau đây đúng hay sai?
%	\choiceTF[t]
{Một vật có khối lượng nhất định nhưng khi ở các thể khác nhau thì sẽ khác nhau về thể tích}
{Một vật có khối lượng nhất định nhưng khi ở các thể khác nhau thì sẽ khác nhau về kích thước nguyên tử cấu tạo nên vật đó}
{Một vật có khối lượng nhất định nhưng khi ở các thể khác nhau thì sẽ khác nhau về khối lượng riêng của chất tạo nên vật đó ở các thể đó}
{Một vật có khối lượng nhất định nhưng khi ở các thể khác nhau thì sẽ khác nhau về trật tự của các nguyên tử cấu tạo nên vật đó}
%	\loigiai{
%		\begin{itemize}
%			\item Đúng
%			\item Sai
%			\item Đúng
%			\item Đúng
%		\end{itemize}
%		Khi chuyển thể, thể tích bị thay đổi, dẫn đến khối lượng riêng thay đổi.
%	}
%\end{ex}
\loai{Loại 15: Giải thích hiện tượng}

\begin{ex} \nguon{4}
	%Dựa vào cấu trúc vật chất theo mô hình động học phân tử hãy giải thích vì sao hơi nước có khối lượng riêng nhỏ hơn nước?
	So sánh giữa nước ở dạng lỏng và nước ở dạng hơi (hơi nước) với cùng khối lượng theo mô hình động học phân tử.
	\choiceTF[t]
{Hơi nước có khối lượng riêng lớn hơn nước lỏng}
{Hơi nước có thể tích lớn hơn nước lỏng}
{Trong hơi nước, các phân tử nước có khoảng cách xa nhau hơn so với khi ở dạng nước lỏng}
{Hơi nước và nước lỏng đều có các phân tử di chuyển hỗn loạn}
	\loigiai{
		\begin{itemize}
			\item SAI ĐÚNG ĐÚNG SAI
			\item Theo mô hình động học phân tử, các chất ở thể khí có khoảng cách giữa các phân tử xa hơn khoảng cách giữa các phân tử khi nó ở thể lỏng. 
			\item Vì vậy với cùng một khối lượng chất như nhau, thể tích khi ở thể khí lớn hơn thể lỏng. 
			\item Điều này làm cho khối lượng riêng của một chất ở thể khí nhỏ hơn khối lượng riêng của chất đó khi ở thể lỏng.
		\end{itemize}
	}
\end{ex}


\begin{ex} \nguon{4}
	Một học sinh đổ $100\ ml$ rượu vào $100\ ml$ nước.
	%, kết quả thu được khoảng $190\ ml$ hỗn hợp rượu với nước. Hãy dùng mô hình động học phân tử giải thích sự biến mất của khoảng $10\ ml$ thể tích các chất.
	\choiceTF[t]
{Kết quả thu được 200 ml hỗn hợp}
{Kết quả thu được lớn hơn 200 ml hỗn hợp}
{Kết quả thu được nhỏ hơn 200 ml hỗn hợp}
{Khi trộn rượu và nước, những phân tử rượu đã xen vào khoảng cách giữa những phân tử nước và ngược lại}
	
	\loigiai{
		SAI SAI ĐÚNG ĐÚNG.\\
		Rượu và nước đều có cấu tạo từ các phân tử. Giữa các phân tử nước và các phân tử rượu đều có khoảng cách. Khi trộn rượu với nước, các phân tử rượu đan xen vào khoảng cách giữa các phân tử nước và ngược lại. Sự đan xen lẫn nhau này làm cho thể tích hỗn hợp bị giảm.
	}
\end{ex}
%\begin{ex}[CD] 
%	Thanh sắt được tạo thành từ các phân tử chuyển động không ngừng nhưng tại sao lại không bị tan rã thành các hạt riêng biệt?
%	
%	\loigiai{
	%		
	%		\begin{itemize}
		%			\item Thanh sắt không được tạo thành từ các phân tử mà là từ các nguyên tử sắt. Trong kim loại, như thanh sắt, nguyên tử không tồn tại dưới dạng phân tử như trong một số chất khác.	
		%			\item Cấu trúc của kim loại được mô tả bởi mô hình lưới kim loại, trong đó các nguyên tử sắt tạo thành một lưới không gian mà mỗi nguyên tử đều kết nối với những nguyên tử xung quanh thông qua liên kết kim loại. Liên kết này là kết quả của sự chia sẻ tự do của các electron dẫn năng (electron tự do) trong lưới kim loại.		
		%			\item Do cấu trúc này, thanh sắt không tồn tại ở dạng phân tử riêng biệt và không bị tan rã thành các hạt riêng lẻ như trong các chất phân tử. Sự liên kết mạnh mẽ giữa các nguyên tử sắt trong lưới kim loại làm cho chúng giữ vững cấu trúc của mình. Nguyên tử sắt không tự do di chuyển và tự tạo thành các đơn vị riêng lẻ mà không cần đến sự liên kết phân tử như trong chất phân tử.
		%		\end{itemize}
	%		
	%	}
%	\haidong{8}\end{ex}

\begin{ex}[KNTT]
Hiện tượng chứng tỏ giữa các phân tử có lực đẩy, lực hút.
%Hãy tìm các hiện tượng thực tế chứng tỏ giữa các phân tử có lực đẩy, lực hút.
\choiceTF[t]
{Cho chất khí nhốt vào một xilanh rồi đẩy pittông nén lại. Ta chỉ có thể nén khối khí đến một thể tích nào đó}
{Khi chất lỏng làm ướt thành bình thì mặt chất lỏng gần thành bình là một mặt lõm (hiện tượng dính ướt)}
{Khi chất lỏng không làm dính ướt thành bình thì mặt chất lỏng gần thành bình là một mặt lồi (hiện tượng không dính ướt)}
{Hiện tượng dâng lên hay hạ xuống của mực chất lỏng ở bên trong các ống có bán kính trong nhỏ so với mực chất lỏng ở ngoài (hiện tượng mao dẫn)}
\loigiai{
	\begin{itemize}
		\item ĐÚNG ĐÚNG ĐÚNG ĐÚNG
		\item Ví dụ: Cho chất khí nhốt vào một xilanh rồi đẩy pittông nén lại. Ta chỉ nén khối khí đến một thể tích nào đó thôi vì khi đó lực đẩy giữa các phân tử là rất lớn, chống lại lực nén của pittông.
	\end{itemize}
}
\haidong{8}\end{ex}
\newpage 
\begin{ex} \nguon{1}
	Nước trong ao, hồ, sông, biển.
	% có chứa không khí không? Tại sao không khí nhẹ hơn nước nhiều mà không bay hết lên trên lại trộn lẫn vào trong nước?
	\choiceTF[t]
{Có chứa không khí}
{Không chứa không khí}
{Không khí nhẹ hơn nước nên bay hết lên trên}
{Phân tử khí xen vào giữa các phân tử nước nhờ lực hút giữa các phân tử nước và phân tử khí}
	\loigiai{
		\item ĐÚNG SAI SAI ĐÚNG
		%Nước trong ao, hồ, sông, biển có chứa không khí vì tính chất của chất khí là chuyển động hỗn độn không ngừng về mọi phía nên sẽ có 1 số phân tử khí chui vào trong nước.
		
		Nước trong ao, hồ, sông, biển thực tế không chỉ có chứa nước mà còn có sự hòa tan của khí oxy và khí nitrogen từ không khí xung quanh. Tuy nhiên, không khí trong nước không phải là khí không khí bên ngoài, mà là sự hòa tan của khí từ môi trường xung quanh vào nước.
		
		
		Nguyên nhân chính là do tính chất hòa tan khí của nước. Khí oxy và khí nitrogen hòa tan trong nước nhờ vào lực hút giữa phân tử nước và các phân tử khí. Quá trình hòa tan này xảy ra do sự chuyển động của phân tử nước và khí, cho phép các phân tử khí xen vào giữa các phân tử nước một cách ngẫu nhiên.
		
		
		Các khí này không bay lên trên mặt nước mà được giữ lại bởi nước chủ yếu là do khối lượng riêng của khí trong nước so với nước rất nhỏ, do đó, chúng không thể tạo ra sự phân tầng mạnh mẽ giống như khi khí nằm trên nước. Ngoài ra, các sự kiện lớn.
		
		
		
		
		
		
	}
\end{ex}

\begin{ex} \nguon{1}
Khi đổ nước sôi vào cốc thủy tinh thì cốc hay bị nứt, còn cốc bằng thạch anh lại không bị nứt. %Giải thích?
\choiceTF[t]
{Do thạch anh và thủy tinh đều là chất rắn kết tinh}
{Do thạch anh và thủy tinh đều là chất rắn vô định hình}
{Do thạch anh là chất rắn vô định hình, còn thủy tinh là chất rắn kết tinh}
{Do thạch anh là chất rắn kết tinh, còn thủy tinh là chất rắn vô định hình}
\loigiai{
	immini[thm]{SAI SAI SAI ĐÚNG.\\
		Thạch anh là chất rắn kết tinh, còn thủy tinh là chất rắn vô định hình, nên thạch anh có liên kết giữa các phân tử chắc chắn hơn nhiều so với thủy tinh
	}{\includegraphics[width=2cm]{images/2024_05_26_010fc85c05a422f9b1ffg-02(1)}}
}
\end{ex}

\begin{ex} \nguon{1}
Khi pha nước chanh, người ta thường làm cho đường tan trong nước rồi mới bỏ đá lạnh vào. 
%Vì sao không bỏ đá lạnh vào trước rồi bỏ đường sau? Giải thích điều này như thế nào?
\choiceTF[t]
{Để làm cho thức uống được ngon hơn}
{Để cho thuận tiện}
{Vì nước ấm thì các phân tử nước chuyển động nhanh hơn, va chạm mạnh hơn với đường khiến đường dễ tan}
{Nếu cho đá vào trước thì nước sẽ lạnh đi khiến cho các phân tử nước chuyển động chậm hơn và va chạm yếu hơn với đường và làm cho đường khó tan hơn}
\loigiai{
	immini[thm]{\begin{itemize}
			\item SAI SAI ĐÚNG ĐÚNG
			\item Hòa tan đường vào nước ấm rồi mới cho đá vào sau thì trong nước ấm các phân tử nước chuyển động nhanh hơn, va chạm mạnh hơn với các phân tử đường nên đường dễ hòa tan.
			\item Còn cho đá vào trước sẽ khiến nước bị lạnh, phân tử nước chuyển động chậm, số lần va chạm với phân tử đường trong một đơn vị thời gian và cường độ va chạm đều giảm nên đường lâu tan hơn.
	\end{itemize}}
	{
\includegraphics[width=5cm]{images/2024_05_26_010fc85c05a422f9b1ffg-02(2)}}
}
\end{ex}





\begin{ex} \nguon{1}
Lấy một cốc nước gần đầy, nếu chỉ cần cho thêm một thìa nước vào trong cốc thì nước sẽ tràn ra nhưng cho thêm một thìa đường vào trong cốc thì nước lại không tràn ra. 
%Hãy giải thích hiện tượng đó?
\choiceTF[t]
{Do giữa các phân tử nước có khoảng cách nên phân tử đường có thể xen vào khoảng trống đó}
{Do đường là chất rắn}
{Do nước có khối lượng riêng nhỏ hơn đường}
{Do đường có khối lượng riêng nhỏ hơn nước}
\loigiai{
	\begin{itemize}
		\item ĐÚNG SAI SAI SAI
		\item Do giữa các phân tử nước có khoảng cách nên các phân tử đường có thể xen vào khoảng trống giữa các phân tử nước và không làm tăng thể tích chiếm chỗ nên nước không tràn ra.
		\item Còn nếu cho thêm thìa nước mới vào thì nước không xen được vào khoảng trống giữa các phân tử nước cũ nên sẽ làm tăng thể tích và nước sẽ tràn ra ngoài.\end{itemize}
}
\end{ex}

%\begin{ex} \nguon{1}
%	Đổ $50 \ cm^3$ rượu vào trong $50 \ cm^3$ nước ta thu được hỗn hợp có thể tích lớn hơn, bằng hay bé hơn $100 \ cm^3$? Giải thích?
%	\begin{itemize}
%		\item[a)] \ds 
%		\item[b)] \ds 
%		\item[c)] \ds 
%		\item[d)]  \ds 
%	\end{itemize}
%	\loigiai{
%		\begin{itemize}
	%			\item Thể tích hỗn hợp sẽ bé hơn $100 \ cm^3$. 
	%			\item Mặc dù cùng là chất lỏng nhưng do kích thước các phân tử và khoảng cách giữa các phân tử rượu khác kích thước phân tử và khoảng cách giữa các phân tử nước nên khi đổ lẫn vào nhau các phân tử rượu có thể xen vào khoảng trống giữa các phân tử nước và làm cho thể tích hỗn hợp giảm đi.\end{itemize}
%	}
%\end{ex}

\begin{ex} \nguon{4}
	Biết rằng các phân tử khí trong phòng chuyển động với tốc độ trung bình từ $100 \ m/s$ đến $2000 \ m/s$. Tuy vậy khi mở lọ nước hoa ở một góc trong phòng thì khoảng vài giây sau những người ở vị trí khác trong phòng mới ngửi thấy mùi nước hoa.  Vì sao có sự trễ như vậy?
	\choiceTF[t]
{Do các phân tử nước hoa không chuyển động thẳng từ lọ nước hoa đến vị trí của những người khác trong phòng}
{Các phân tử nước hoa chuyển động hỗn loạn dích dắc từng đoạn rất ngắn do bị va chạm với các phân tử không khí}
{Do gió trong phòng nên các phân tử nước hoa không đi thẳng được}
{Do trong nước hoa có chứa cồn}
	\loigiai{
		\begin{itemize}
			\item  ĐÚNG ĐÚNG SAI SAI
			\item Do các phân tử nước hoa không chuyển động thẳng từ lọ nước hoa đến vị trí của những người khác trong phòng, mà đó là chuyển động hỗn loạn dích dắc từng đoạn rất ngắn do bị va chạm với các phân tử không khí. 
			\item Vì vậy thời gian chuyển động của các phân tử nước hoa lâu hơn.
		\end{itemize}
	}
\end{ex}


\begin{ex} \nguon{4}
Một học sinh nhỏ vài giọt thuốc tím vào hai cốc nước, một cốc nước nóng và một cốc nước lạnh. Quan sát sự khuếch tán của thuốc tím ở hai cốc.%Sự khuếch tán của thuốc tím ở cốc nào diễn ra nhanh hơn? Giải thích.
\choiceTF[t]
{Cốc nước nóng diễn ra nhanh hơn}
{Cốc nước nóng diễn ra chậm hơn}
{Hai cốc diễn ra bằng nhau}
{Các phân tử ở cốc nước nóng chuyển động nhanh hơn tác động vào phân tử thuốc tím chuyển động hỗn loạn hơn}
\loigiai{
	ĐÚNG SAI SAI ĐÚNG
	Sự khuếch tán của thuốc tím ở cốc nước nóng diễn ra nhanh hơn.
	Giải thích: cốc nước nóng có nhiệt độ cao làm cho sự chuyển động hỗn loạn của các phân tử càng nhanh hơn.
}
\end{ex}

\begin{ex} \nguon{4}
%Dựa vào cấu trúc và sự sắp xếp của các phân tử trong chất lỏng và chất khí. 
%Hãy giải thích vì sao trong chất lỏng sự khuếch tán diễn ra chậm hơn rất nhiều lần so với trong chất khí?
So sánh sự khuếch tán trong chất lỏng và chất khí.
\choiceTF[t]
{Chất lỏng diễn ra chậm hơn vì mật độ phân tử lớn  hơn}
{Chất khí diễn ra chậm hơn vì mật độ phân tử nhỏ hơn}
{Chất lỏng diễn ra chậm hơn vì liên kết giữa các phân tử chặt chẽ hơn}
{Cả hai diễn ra như nhau}
\loigiai{
	ĐÚNG SAI ĐÚNG SAI.\\
	Mật độ phân tử ở trạng thái lỏng lớn hơn rất nhiều (khoảng 1000 lần) mật độ phân tử ở trạng thái khí. Sự khuếch tán ra xa của một phân tử nào đó trong chất lỏng diễn ra chậm chạp hơn là vì phân tử này va chạm nhiều lần hơn với các phân tử khác so với khi nó khuếch tán trong chất khí. Mặt khác sự liên kết giữa các phân tử chất lỏng chặt chẽ hơn các phân tử chất khí, điều này cũng cản trở sự khuếch tán.
}
\end{ex}

\begin{ex} \nguon{4}
%Giải thích vì sao 
Khi trời lạnh, quá trình bôi trơn của dầu nhớt trong động cơ ô tô, xe máy... hoạt động kém hiệu quả.
\choiceTF[t]
{Dầu nhớt thường được tạo thành từ các phân tử dài, có cấu trúc phức tạp}
{Do những phân tử dầu nhớt thường có khả năng trượt qua nhau một cách dễ dàng khi nhiệt độ tăng lên}
{Do những phân tử dầu nhớt thường có khả năng trượt qua nhau một cách dễ dàng khi nhiệt độ giảm xuống}
{Do khi đó dầu nhớt loãng hơn}
\loigiai{
	ĐÚNG ĐÚNG SAI SAI.\\
	Dầu nhớt thường được tạo thành từ các phân tử dài, có cấu trúc phức tạp. Những phân tử này thường có khả năng trượt qua nhau một cách dễ dàng khi nhiệt độ tăng lên. Tuy nhiên, khi nhiệt độ giảm, các phân tử này sẽ có xu hướng tự sắp xếp thành các cấu trúc có trật tự hơn, không còn chuyển động hỗn loạn mạnh, làm tăng độ đặc của dầu. Điều này làm giảm hiệu quả bôi trơn của dầu nhớt.
}
\end{ex}

\begin{ex} \nguon{4}
Khi tiến hành quét sơn lên bề mặt vật liệu.
%Hãy dùng thuyết động học phân tử về sự tương tác giữa các phân tử để giải thích sự kết dính của các phân tử nước sơn và bề mặt vật liệu.
\choiceTF[t]
{Các phân tử sơn tạo ra một lực hút với bề mặt vật liệu}
{Một số phân tử sơn khác có thể tạo ra một lực đẩy để giữ cho sơn trải đều trên bề mặt}
{Sự kết dính được tạo ra bởi các lực tương tác giữa các phân tử sơn và bề mặt của vật liệu}
{Các phân tử sơn xen vào sâu bên trong vật liệu}
\loigiai{
	ĐÚNG ĐÚNG ĐÚNG SAI\\
	Khi tiến hành sơn các bề mặt vật liệu, sự kết dính được tạo ra bởi các lực tương tác giữa các phân tử sơn và bề mặt của vật liệu mà sơn được sử dụng để dính. Các phân tử sơn tạo ra một lực hút với bề mặt vật liệu. Trong khi đó một số phân tử sơn khác có thể tạo ra một lực đẩy để giữ cho sơn trải đều trên bề mặt.
}
\end{ex}
\begin{ex} \nguon{4}
Cao su là một loại vật liệu có tính đàn hồi. 
%Hãy dùng thuyết động học phân tử giải thích tính đàn hồi của cao su khi nó bị kéo giãn hoặc khi bị nén.
\choiceTF[t]
{Các chuỗi phân tử cao su thường tạo ra các tương tác hút lẫn nhau thông qua các lực hút}
{Trong cấu trúc của cao su, các phân tử tạo ra lực đẩy lẫn nhau, đặc biệt là khi cao su bị nén hoặc uốn cong}
{Cao su là chất rắn kết tinh}
{Tính đàn hồi của cao su khi nó bị kéo dãn hoặc khi bị nén là do lực hút và lực đẩy phân tử}
\loigiai{
	ĐÚNG ĐÚNG SAI ĐÚNG\\
	Các chuỗi phân tử cao su thường tạo ra các tương tác hút lẫn nhau thông qua các lực hút. Các tương tác này giữ cho các chuỗi phân tử gần nhau và giữ cho cao su giữ được hình dạng và kết cấu ban đầu. Ngoài ra, trong cấu trúc của cao su, các phân tử cũng tạo ra lực đẩy lẫn nhau, đặc biệt là khi cao su bị nén hoặc uốn cong. Sự đẩy này giúp duy trì tính đàn hồi và co giãn của vật liệu khi tác dụng lực lên nó.
}
\end{ex}

\begin{ex} \nguon{4}
Biết rằng muối ăn là chất rắn kết tinh, nhựa thông là chất rắn vô định hình. Ta đập vỡ vụn hạt muối ăn và cục nhựa thông rồi quan sát các mảnh vụn đó dưới kính lúp.
\choiceTF[t]
{Vụn muối ăn vẫn có dạng hình học xác định}
{Vụn nhựa thông có hình dạng bất kì}
{Vụn muối ăn vẫn có hình dạng bất kì}
{Cả hai đều có hình dạng bất kì}
\loigiai{
	ĐÚNG ĐÚNG SAI SAI\\
	Vụn muối ăn vẫn có dạng hình học. Vì muối ăn là chất rắn kết tinh, có cấu trúc mạng tinh thể xác định, khi bị vỡ vẫn còn giữ được cấu trúc mạng tinh thể.
	Vụn nhựa thông có hình dạng bất kì. Vì nhựa thông là chất rắn vô định hình, không có cấu trúc tinh thể, khi vỡ các mảnh vụn sẽ theo các hình dạng bất kì, không xác định.
}
\end{ex} 
\begin{ex} \nguon{4} 
Kim cương và than chì đều có cấu trúc dưới dạng chất rắn, cùng được cấu tạo từ nguyên tố carbon nhưng lại có tính chất khác nhau, điển hình kim cương rất cứng còn than chì thì mềm. %Nguyên nhân nào đã dẫn đến sự khác nhau như vậy?
\choiceTF[t]
{Sự khác nhau về tính chất của kim cương và than chì là do cấu trúc mạng tinh thể của chúng}
{Mỗi nguyên tử carbon trong tinh thể kim cương liên kết với bốn nguyên tử khác tạo nên một cấu trúc tinh thể vững chắc}
{Than chì có cấu trúc tinh thể dạng xếp thành lớp, mỗi nguyên tử carbon liên kết với ba nguyên tử carbon khác. Cấu trúc này không tạo nên sự bền vững cho than chì}
{Do kim cương là chất rắn kết tinh còn than chì là chất rắn vô định hình}
\loigiai{
	immini[thm]{ĐÚNG ĐÚNG ĐÚNG SAI\\
		Sự khác nhau về tính chất của kim cương và than chì là do cấu trúc mạng tinh thể của chúng.
		
		
		Mỗi nguyên tử carbon trong tinh thể kim cương liên kết với bốn nguyên tử khác tạo nên một cấu trúc tinh thể vững chắc}{\includegraphics[width=3cm]{images/2024_05_29_5f55f4890f9e0910bb1cg-06(1)}}
	immini[thm]{Than chì có cấu trúc tinh thể dạng xếp thành lớp, mỗi nguyên tử carbon liên kết với ba nguyên tử carbon khác. Cấu trúc này không tạo nên sự bền vững cho than chì}{\includegraphics[width=3cm]{images/2024_05_29_5f55f4890f9e0910bb1cg-06}}
}
\end{ex}

\begin{ex} \nguon{4}
Thanh kim loại ở nhiệt độ phòng thì rất cứng nhưng khi nung nóng lại trở nên mềm dẻo, dễ uốn cong hơn. Nguyên nhân là vì
\choiceTF[t]
{khi được nung nóng thì kim loại biến thành chất rắn vô định hình}
{khi được nung nóng thì cấu trúc của kim loại trở nên linh hoạt hơn và dẻo hơn}
{khi được nung nóng các phân tử trong kim loại có khả năng trượt qua nhau dễ dàng khiến cho kim loại trở nên mềm dẻo hơn}
{khi nhiệt độ tăng lên, dao động nhiệt của các phân tử trong kim loại cũng tăng khiến cho liên kết giữa các phần tử trở nên yếu dần}
\loigiai{
	SAI ĐÚNG ĐÚNG ĐÚNG.\\
	Khi nhiệt độ tăng lên, dao động nhiệt của các phân tử trong kim loại cũng tăng khiến cho liên kết giữa các phần tử trở nên yếu dần. Kết quả là cấu trúc của kim loại trở nên linh hoạt hơn và dẻo hơn.
	Do vậy, khi được nung nóng các phân tử trong kim loại có khả năng trượt qua nhau dễ dàng khiến cho kim loại trở nên mềm dẻo hơn. Lúc này thanh kim loại dễ uốn cong hoặc định hình lại một cách dễ dàng.
}
\end{ex}

%\begin{ex} \nguon{4}
%	Hai vật $A$ và $B$ giống nhau nhưng vật $A$ có nhiệt độ cao hơn vật $B$. Hãy cho biết sự khác nhau về trạng thái chuyển động của các phân tử cấu tạo của hai vật $A$ và $B$.
%	\begin{itemize}
%		\item[a)] \ds 
%		\item[b)] \ds 
%		\item[c)] \ds 
%		\item[d)]  \ds 
%	\end{itemize}
%	\loigiai{
%		Vật $A$ có nhiệt độ cao hơn vật $B$. Vì vậy các phân tử cấu tạo nên vật $A$ có chuyển động nhiệt nhanh hơn các phân tử cấu tạo nên vật $B$.
%	}
%\end{ex}
\begin{ex}
	immini[thm]{Thí nghiệm nén khí và nén nước\\
		\textit{Dụng cụ:} xi-lanh, pít-tông, nước, nút bấc.\\
		\textit{Tiến hành}: Kéo pít-tông để hút một lượng không khí vào xi-lanh. Dùng nút bấc nút chặt đầu xi-lanh rồi ấn pít-tông để nén không khí.
		Sau đó, cho một lượng nước vào xi-lanh và lặp lại thao tác như trên để nén nước trong xi-lanh. Trường hợp nào nén dễ dàng hơn? Hãy giải thích hiện tượng bằng mô hình động học phân tử}{\vspace*{-.5 cm} includegraphics[width=8cm]{images/xilanhnew}}
	
	\loigiai{
		\begin{itemize}
			\item Nén không khí dễ dàng hơn so với nén nước. Điều này được giải thích bằng mô hình động học phân tử như sau:
			\begin{itemize}
				\item Không khí là một chất khí, các phân tử ở xa nhau và chuyển động tự do, do đó có nhiều khoảng trống giữa các phân tử và dễ bị nén.
				\item 		 Nước là một chất lỏng, các phân tử ở gần nhau hơn và có lực hút mạnh hơn giữa các phân tử, do đó khó bị nén hơn nhiều so với không khí.
			\end{itemize}
		\end{itemize}
	}
\end{ex}
